\chapter{Varietà differenziabili}\stepcounter{capitoloeccolo}

\section{Prevettori}

\definizione{
    Per $ r, s \in \starR $, si adotterà la seguente notazione: \begin{itemize}
        \item $ r \prec s $ se $ r = a \cdot s $ per qualche $ a $ \emph{finito};
        \item $ r \pprec s $ se $ r = \lambda \cdot s $ per qualche $ \lambda $ infinitesimo.
    \end{itemize}
}
\osservazione{
    Questa definizione si estende naturalmente ad elementi di $ \nonstandard{(\R^n)} $. In particolare, se $ v \in \nonstandard{(\R^n)} $ e $ s \in \starR $, si adotta la seguente notazione: \begin{itemize}
        \item $ v \prec s $ se $ \nonstandard{\norm{v}} \prec s $ per qualche norma $ \norm{ \cdot } $ su $ \R^{n} $;
        \item $ v \pprec s $ se $ \nonstandard{\norm{v}} \pprec s $ per qualche norma $ \norm{ \cdot } $ su $ \R^{n} $;
    \end{itemize}

    Se $ v, w \in \nonstandard{(\R^n)} $, si scriverà $ v\approx w $ se $ v-w \pprec 1 $.

    In particolare, se $ v = (v_1,\dots,v_{n} ) \in \nonstandard{(\R^n)} $, si ha che $ v \prec s $ oppure $ v \pprec s $ se e solo se ogni $ v_{i}  $ soddisfa $ v_{i} \prec s  $ oppure $ v_{i} \pprec s $. 
}
\osservazione{
    Poiché su $ \R^{n} $ tutte le norme sono equivalenti, la definizione precedente è ben posta.
} %% Osservazione: in uno spazio vettoriale finito-dimensionale tutte le norme sono equivalenti

\paragrafo{Osservazione}{%
    Per ogni $ a,b \in \starR $, vale almeno una delle seguenti: \begin{itemize}
        \item $ a \prec b $; 
        \item $ b\pprec a $.
    \end{itemize}Infatti, se per assurdo esistessero $ a,b \in \starR $ tali per cui nessuna delle due fossero verificate, allora \begin{itemize}
        \item $ a/b $ è infinito,
        \item $ b/a $ non è infinitesimo,
    \end{itemize}e questo genera un assurdo.
}{djknckjnsjdjdjcskjnskjnckjnskjnd}{}

\definizione{
    Sia $ X $ uno spazio topologico e sia $\nonstandard{X} $ la sua estensione non standard.\begin{itemize}
        \item $ \forall\, p \in X $ si definisce l'\emph{alone} di $ p $: \[
            \halo(p) = \bigcap_{U \in I(p) }\hspace{-.5em}\nonstandard{U} 
        \]dove $ I(p) $ è l'insieme di tutti gli intorni aperti di $ p $.
        \item I punti di $ X $ si chiamano \emph{standard}, e un punto che fa parte di $ \halo(p) $ per qualche $ p \in X$ è detto \emph{quasi-standard}: \begin{equation*}
            \hhalo{X} \coloneqq \bigcup_{p \in X}\halo(p) .
        \end{equation*}
        \item Se $ a $ è un punto quasi-standard allora la \emph{parte standard} di $ a $, denotata con $ \st(a)$, è l'unico $ p \in X$ tale che \[
            a \in \halo(p).
        \]
        \item Per un generico $ A \subseteq X $ si definisce l'\emph{alone} di $ A $ come \[
            \hhalo{A} = \bigcup_{a \in A} \halo(a).
        \]
    \end{itemize} 
}
\osservazione{
    Se $ M $ è una varietà differenziabile, e $ p \in M $, $ \halo(p) \subseteq \nonstandard{M} $ è l'insieme dei punti $ x \in \nonstandard{M} $ per i quali vi è un intorno coordinato $ U $ di $ p $ tale che $ x \in \nonstandard{U} $, e $ x \approx p $ in quelle date coordinate.
}
\proprieta{
    Per uno spazio topologico $ X $ (con estensione non standard $ \nonstandard{X} $) valgono le seguenti proprietà. \begin{enumerate}
        \item Se $ A \subseteq X $ è aperto in $ X $, allora $ \hhalo{A} \subseteq \nonstandard{A} $. 
        \item Se $ A \subseteq X $ è compatto in $ X $, allora $ \nonstandard{A} \subseteq \hhalo{A} $
        \item Se $ X $ è compatto allora $\hhalo{X} = \nonstandard{X}$
    \end{enumerate}
}
\lemma{kjnsckjndjcjsjnckjnsdkjncskjndkjncskjnkjnj}{
    Per un aperto $ W \subseteq  \R^{n} $, sia $ f: W \longrightarrow \R $ continua, e sia \[
        \nonstandard{f}: \nonstandard{W} \longrightarrow \nonstandard{\R}
    \]la sua estensione non standard. Allora $ \nonstandard{f}(a)$ è finita per ogni $ a \in \hhalo{W} $.
}
\dimostrazionelem{kjnsckjndjcjsjnckjnsdkjncskjndkjncskjnkjnj}{
    Sia $ a \in \null^{\halo}W $, e sia $ U $ un intorno di $ \st(a) $ tale che $ \overline{U} \subseteq W $ e $ \overline{U} $ sia compatto. 

    Dunque $ \exists\, C \in \R $ tale che $ \forall\, x \in U $: $ |f(x)| \le C $.
        
    Per il principio di transfer, $\forall\, x \in \nonstandard{U}$: $ |\nonstandard{f}(x)| \le C$. In particolare, $ a \in \nonstandard{U} $, e dunque \[
        |\nonstandard{f}(a)| \le C.\qedd
    \]
}
\paragrafo{Osservazione}{%
    Dato un aperto $ U \subseteq \R^{n} $ ed una funzione $ f: U \longrightarrow \R $, le derivate parziali di $ \nonstandard{f} $, quando esistono, sono definite come l'estensione non standard delle derivate parziali di $ f $: \[
        \nonstandard{\left(\dpd{f}{x_{i} }\right)}: \nonstandard{U} \longrightarrow \nonstandard{\R}.
    \]\begin{itemize}
        \item Se $ f \in C^{1}(U) $, allora per ogni $ a \in \nonstandard{U} $ esiste il vettore $ D_{a} $ di derivate parziali nel punto $ a $.
        \item Se $ f \in C^{2}(U) $, allora per ogni $ a \in \nonstandard{U} $ esiste la matrice Hessiana $ H_{a}$ di derivate seconde nel punto $ a $.
    \end{itemize}
}{}{}
\paragrafo{Proprietà}{%
    Sia $ U \subseteq \R^{n} $ aperto e sia $ f : U \longrightarrow \R $. Siano $ a,b \in \hhalo{U} $ con $ a \approx b $: allora l'intervallo $ [a,b] \subseteq \hhalo{U} \subseteq \nonstandard{U}. $
    
    \begin{enumerate}
        \item  Se $ f \in C^{1} $, allora per il transfer del teorema del valor medio, esiste $ x \in [a,b] $ tale che \begin{equation}
            \nonstandard{f}(b)-\nonstandard{f}(a) = D_{x} (b-a).\label{eq:1}
    \end{equation}Poiché le derivate parziali sono continue, si ha che\[
            \forall\, y \approx a:\qquad D_{x}-D_{y} \pprec 1 
        \]dunque, manipolando \eqref{eq:1}: \[
            \nonstandard{f}(b)-\nonstandard{f}(a) = D_{y}(b-a) + (D_{x}-D_{y}  ) (b-a) 
        \] 
        
        $\implies$ $ \nonstandard{f}(b)-\nonstandard{f}(a) - D_{y}(b-a) \pprec \norm{b-a} $.

        Se inoltre le derivate parziali sono Lipschitziane ed esiste $ \beta \pprec 1 $\footnote{cioè $ \beta $ infinitesimo} tale che \[
            \forall\, \xi \in [a,b]\qquad \norm{\xi-y}\prec \beta
        \]allora vale la condizione più forte \[
            \nonstandard{f}(b)-\nonstandard{f}(a) - D_{y}(b-a) \prec \beta\norm{b-a}.
        \]
        \item Se $ f \in C^{2} $, allora per il transfer del teorema di approssimazione di Taylor si ha \[
            \nonstandard{f}(b)-\nonstandard{f}(a) = D_{a} (b-a) + \frac{1}{2} \null^{t}\!(b-a) H_{x} (b-a) 
        \]per qualche $ x \in [a,b] $.
        \item La proprietà 1 si può applicare a ciascuna delle componenti di $ \varphi $: \begin{align*}
            \varphi: U & \longrightarrow \R^{n}\\
            x &\longmapsto \left(\varphi_{1}(x), \dots, \varphi_{n}(x)  \right)
        \end{align*}ottenendo così una condizione sulla matrice Jacobiana $ J_{a}  $ (composta dalle $ n $ righe $ D_{a}^{i}  $ corrispondenti alle $ \varphi_{i}  $): \[
            \nonstandard{\varphi}(b)-\nonstandard{\varphi}(a)-J_{y}(b-a)\pprec \norm{b-a}\qquad\forall y \approx a 
        \]

        Se inoltre tutte le derivate parziali delle $ \varphi_{i}  $ sono Lipschitziane ed esiste $ \beta \pprec 1 $ tale che \[
            \forall\, \xi \in [a,b]\qquad \norm{\xi-y}\prec \beta
        \]allora vale la condizione più forte \[
            \nonstandard{\varphi}(b)-\nonstandard{\varphi}(a) - J_{y}(b-a) \prec \beta\norm{b-a}.
        \]
    \end{enumerate}
}{dckjndjjjjjjdjscacchina}{}
\paragrafo{Notazione}{%
    Si fissano ora delle convenzioni valide per tutto il capitolo. \begin{itemize}
        \item Sia $ M $ la varietà differenziabile di dimensione $ n $ di interesse, definita come nella Sezione \ref{sec:geodiff}. 
        \item D'ora in avanti sia $ \lambda \in \starR $ un infinitesimo positivo \emph{fissato}.
    \end{itemize}
}{}{}


\definizione{
    Per $ a \in \hhalo{M} $, un \emph{prevettore} in $ a $ è una coppia $ (a,x) $ con $ x \in \hhalo{M} $ tale che per ogni funzione differenziabile $ f: M \longrightarrow \R $, \[
        \nonstandard{f}(x)-\nonstandard{f}(a)\prec \lambda
    \]Si indicherà con $ P_{a}M  $ l'insieme dei prevettori in $ a $.
}
\proposizione{dkjnckjnsckjnsjdjckjnsjcjdkjn}{
    Sia $ a \in \hhalo{M} $. Sono fatti equivalenti: \begin{enumerate}
        \item $ (a,x) $ è un prevettore in $ a $. 
        \item Se $ W $ è un intorno coordinato di $ \st(a) $ la cui immagine è $ U \subseteq \R^{n} $, e $ \hat{a}, \hat{x} \in \nonstandard{U} $ sono le coordinate di $ a $ e $ x $ rispettivamente: \[
            \hat{x}-\hat{a}\prec \lambda.
        \]
    \end{enumerate}
}
\dimostrazioneprop{dkjnckjnsckjnsjdjckjnsjcjdkjn}{
    $\boxed{1.\implies 2.}$\hspace{1em} Siano $ x_1,\dots,x_{n} $ le funzioni coordinate di $ W $: poiché ciascuna $ x_{i} $ è differenziabile, per ipotesi \[
        \nonstandard{x_{i} }(x)-\nonstandard{x_{i} }(a)\prec \lambda  
    \]e dunque $ \hat{x}-\hat{a}\prec \lambda $.

    $\boxed{2.\implies 1.}$\hspace{1em} Sia $ f: M \longrightarrow \R $ una generica funzione differenziabile. Allora esiste $ c \in [\hat{a}, \hat{x}] $ tale che\[
        f(x)-f(a) = D_{c} (\hat{x}-\hat{a}) 
    \]dove le componenti di $ D_{c}$ sono finite per il Lemma \framref{kjnsckjndjcjsjnckjnsdkjncskjndkjncskjnkjnj}. Si può dunque concludere che  \[f(x)-f(a)\prec\norm{\hat{x}-\hat{a}}\prec \lambda \qedd\]
}
\definizione{
    Si definisce una relazione di equivalenza $ \equiv $ in $ P_{a}M  $ come segue: \[
        (a,x) \equiv (a,y) \qquad \text{se e solo se}\qquad \nonstandard{f}(y)-\nonstandard{f}(x)\pprec \lambda
    \]per ogni funzione differenziabile $ f: M \longrightarrow \R $.

    L'insieme delle classi di equivalenza $ P_{a}M/\equiv  $ si denoterà con $ T_{a}(M)  $.
}
\notazione{
    Gli elementi di $ T_{a}M  $, ovvero le classi di equivalenza $ [(a,x)] $ si denoteranno con $ \vec{ax} $.
}
\proposizione{dinckjndjcjdkckjncskjdkckjnk}{
    Sia $ a \in \hhalo{M} $, e siano $ x,y \in \hhalo{M} $. Sono fatti equivalenti: \begin{enumerate}
        \item $ (a,x) \equiv (a,y) $;
        \item Se $ W $ è un intorno coordinato di $ \st(a) $ la cui immagine è $ U \subseteq \R^{n} $, e $ \hat{x}, \hat{y} \in \nonstandard{U} $ sono le coordinate di $ x $ e $ y $ rispettivamente: \[
            \hat{y}-\hat{x}\pprec \lambda.
        \]
    \end{enumerate}
}
\proposizione{dkjnckjnskjndjcjskjncjsjdkjncjdkjkjncjsjdkcjdjckjnscdd}{
    $ T_{a}M  $ è uno spazio vettoriale \emph{reale} $ n $-dimensionale indipendente dalla scelta delle coordinate.
}
\dimostrazioneprop{dkjnckjnskjndjcjskjncjsjdkjncjdkjkjncjsjdkcjdjckjnscdd}{
    Sia $ a \in \hhalo{M} $ fissato, $ W \subseteq M $ intorno coordinato di $ \st(a) $. Si ha quindi una mappa coordinata \begin{align*}
    W &\longrightarrow U \subseteq \R^{n} \\
    x &\longmapsto \hat{x}
    \end{align*}

    Si definisce ora \begin{align*}
        (\R^{n})^{\lambda}_{\text{F}} &\coloneqq \set{v \in \nonstandard{(\R^n)}: v\prec \lambda}\\ 
        (\R^{n})^{\lambda}_{\text{I}} &\coloneqq \set{v \in \nonstandard{(\R^n)}: v\pprec \lambda}
    \end{align*}Il Teorema 10.3.8 di \cite{katz_true_nodate} garantisce che $ (\R^{n})^{\lambda}_{\text{F}}/(\R^{n})^{\lambda}_{\text{I}} \cong \R^{n}$. Dunque, la mappa suriettiva \begin{align*}
    P_{a} &\longrightarrow (\R^{n})^{\lambda}_{\text{F}} \\
    (a,x) &\longmapsto \hat{x}-\hat{a}
    \end{align*}induce una mappa iniettiva \begin{align*}
        T_{a} = \frac{P_{a}}{\equiv}\hspace{0.3em} &\hookrightarrow\hspace{0.3em}  \frac{(\R^{n})^{\lambda}_{\text{F}}}{(\R^{n})^{\lambda}_{\text{I}}} \cong \R^{n}\\
        \vec{ax}\hspace{0.3em} &\mapsto \hspace{0.3em} \hat{x}-\hat{a} + (\R^{n})^{\lambda}_{\text{I}}
    \end{align*}che è quindi biiettiva. Questo induce una struttura di spazio vettoriale reale su $ T_{a}$.

    Inoltre, se $ W' $ è un differente intorno coordinato di $ \st(a) $, con immagine $ U' $, e $ \varphi: U \longrightarrow U' $ è il cambio di coordinate, allora per il punto 3. della Proprietà \framref{dckjndjjjjjjdjscacchina} si ha \begin{equation*}
        \varphi(\hat{x})-\varphi(\hat{a}) - J_{\st(a)} (\hat{x}-\hat{a}) \pprec \norma{\hat{x}-\hat{a}} \prec \lambda
    \end{equation*}Dunque la mappa $ \R^{n} \longrightarrow \R^{n} $ indotta dalle due identificazioni di $ T_{a} $ con $ \R^{n} $  fornite dalle due scelte di coordinate è data dalla moltiplicazione per una matrice $ J_{\st(a)} $, e pertanto è lineare. Dunque la struttura di spazio vettoriale di $ T_{a}$ su $ \R $ è ben definita.\qed
}
\osservazione{
    Se $ a,b \in \hhalo{M} $ e $ a \approx b $, allora date le coordinate in un intorno di $ \st(a) $, l'identificazione di $ T_{a}M  $ e $ T_{b}M  $ con $ \R^{n}$ indotta dalle coordinate induce una identificazione tra $ T_{a}M  $ e $ T_{b}M  $. Per una diversa scelta di coordinate, la matrice jacobiana $ J_{\st(a)}  $ definita nella dimostrazione della Proposizione \framref{dkjnckjnskjndjcjskjncjsjdkjncjdkjkjncjsjdkcjdjckjnscdd} è la stessa per $ a $ e per $ b $, e dunque l'identificazione tra $ T_{a}M$ e $ T_{b}M  $ è indipendente dalle coordinate, e dunque ben definita. 

    Di conseguenza, quando $ a \approx b $ è possibile sommare senza ambiguità due vettori $ \vec{ax} \in T_{a}M  $ e $ \vec{by} \in T_{b}M  $.
}
\definizione{
    Sia $ a \in \hhalo{M} $. Un prevettore $ (a,x) \in P_{a}M  $ agisce su una funzione differenziabile $ f: M \longrightarrow \R $ come segue: \[
        (a,x)\,f \coloneqq \frac{1}{\lambda}\left(\nonstandard{f}(x)-\nonstandard{f}(a)\right) \in \starR
    \]che è finito per definizione di prevettore.
}
\definizione{
    Sia $ a \in \hhalo{M} $. Un vettore $ \vec{ax} \in T_{a}M  $ agisce come derivazione su una funzione differenziabile $ f: M \longrightarrow \R $ come segue: \[
        \vec{ax} \,f = \st\left((a,x)\,f\right) \in \R.
    \] 
}

    La definizione di derivazione è ben definita per la definizione della classe di equivalenza $ \equiv $ su $ P_{a}M  $.

\proposizione{dkjncjdjckjnskjnckjdjckjnkjnkjnkjn}{
    L'azione di un prevettore rispetta la regola di Leibniz fino agli infinitesimi, ovvero, se $ a \in \hhalo{M} $, $ (a,x) \in P_{a}M  $ e $ f,g: M \longrightarrow \R $ sono funzioni differenziabili, allora \begin{equation}
        (a,x)\,(fg) \approx \nonstandard{f}(a) \cdot (a,x)\, g + (a,x)\,f\cdot \nonstandard{g}(a).\label{eq:eqdddd}
    \end{equation}

    Inoltre, se $ a_0 = \st(a) $, si ha \begin{equation}
        \vec{ax}\,(fg) = f(a_0) \cdot \vec{ax}\,g + \vec{ax}\,f \cdot g(a_0).\label{eq:djcjnnnnnn}
    \end{equation}
}
\dimframmento{dkjncjdjckjnskjnckjdjckjnkjnkjnkjn}{La \eqref{eq:eqdddd} si dimostra come segue:
    \begin{align*}
        (a,x)\,(fg) &= \frac{1}{\lambda}\left(\nonstandard{f}(x)\nonstandard{g}(x)-\nonstandard{f}(a)\nonstandard{g}(a)\right)\\ 
        &= \frac{1}{\lambda}\left(\nonstandard{f}(x)\nonstandard{g}(x)-\nonstandard{f}(x)\nonstandard{g}(a)+\nonstandard{f}(x)\nonstandard{g}(a)-\nonstandard{f}(a)\nonstandard{g}(a)\right)\\  
        &= \nonstandard{f}(x) \cdot \frac{1}{\lambda}\left(\nonstandard{g}(x)-\nonstandard{g}(a)\right) + \nonstandard{g}(a) \cdot \frac{1}{\lambda} \left(\nonstandard{f}(x) - \nonstandard{f}(a)\right)\\ 
        &\approx \nonstandard{f}(a) \cdot \frac{1}{\lambda}\left(\nonstandard{g}(x)-\nonstandard{g}(a)\right) + \nonstandard{g}(a) \cdot \frac{1}{\lambda} \left(\nonstandard{f}(x) - \nonstandard{f}(a)\right)\\ 
        &= \nonstandard{f}(a) \cdot (a,x)\, g + (a,x)\,f\cdot \nonstandard{g}(a).
    \end{align*}

    La \eqref{eq:djcjnnnnnn} è invece una applicazione della \eqref{eq:eqdddd}: \begin{align*}
        \vec{ax}\,(fg) &= \st\left((a,x)\,fg\right)\\ 
        &= \st\left(\nonstandard{f}(a) \cdot (a,x)\, g + (a,x)\,f\cdot \nonstandard{g}(a)\right)\\ 
        &=  \st\left(\nonstandard{f}(a)\right) \cdot \vec{ax}\,g + \st\left(\nonstandard{g}(a)\right) \cdot \vec{ax}\,f\\ 
        &= f(a_0) \cdot \vec{ax}\,g + g(a_0) \cdot \vec{ax}\,f
    \end{align*}dove l'ultima uguaglianza è conseguenza della continuità di $ f $ e $ g $.\qed
}
\teorema{dkjnckakjndjcjskjndjcjskjndjdjckjnsdjkjncjdjdjsk}{
    Ogni vettore classico $ X $ tangente ad $ M $ in $ a \in M$ definisce un unico vettore $ \vec{ax} \in T_{a} M  $ dove, in coordinate, $ x = a + \lambda X $.
}
\dimostrazione{dkjnckakjndjcjskjndjcjskjndjdjckjnsdjkjncjdjdjsk}{
    Sia $ f: M \longrightarrow \R $ una funzione differenziabile. La derivata classica $ Xf $ in $ a_0 \in M$, per definizione, è \begin{equation*}
        Xf = \st\left(\frac{f(a_0+\lambda X)- f(a_0)}{\lambda}\right)
    \end{equation*}e si ha, per il punto 1. della Proprietà \framref{dckjndjjjjjjdjscacchina}, che per $ a_0\approx a $:\begin{equation}
        \frac{f(a_0+\lambda X)- f(a_0)}{\lambda} \approx \frac{f(a+\lambda X)- f(a)}{\lambda}\label{dkjnckjn} 
    \end{equation}Pertanto, poiché per definizione \begin{equation*}
        \vec{ax} f = \st\left((a,x) f\right) = \st\left(\frac{f(a+\lambda X)- f(a)}{\lambda}\right)
    \end{equation*}applicando \eqref{dkjnckjn} si ottiene che $ \vec{ax} $ e $ X $ agiscono su ogni funzione allo stesso modo.\qed
}
\section{Campi di prevettori}

\definizione{%
    Si definisce un \emph{campo di prevettori} su $ \nonstandard{M} $ una mappa \emph{interna} $ F: \nonstandard{M} \longrightarrow \nonstandard{M} $ tale che $ \left(a, F(a)\right) \in P_{a}M  $ per ogni $ a \in \hhalo{M} $.
} 

    La definizione precedente è equivalente a richiedere che \begin{equation*}
        F(a)-a \prec \lambda\qquad \forall\, a \in \hhalo{M}.
    \end{equation*}

\definizione{
    Se $ F,G: \nonstandard{M} \longrightarrow \nonstandard{M} $ sono due campi di prevettori, $ F $ e $ G $ si dicono \emph{equivalenti} (e si scrive $ F\equiv G $) se \begin{equation*}
        \left(a,F(a)\right) \equiv \left(a,G(a)\right)\qquad \forall\,a \in \hhalo{M}
    \end{equation*}
}
\definizione{
    Un campo di prevettori \emph{locale} è una mappa interna \begin{equation*}
        F: \nonstandard{U} \longrightarrow \nonstandard{V}
    \end{equation*}che soddisfi la condizione di cui sopra, per $ U \subseteq V \subseteq M $ aperti.
}

    Segue naturalmente che un campo di prevettori definito su tutto $ \nonstandard{M} $ viene detto \emph{campo di prevettori globale}. 

\definizionegenerale{Definizione}{
    Un campo locale di prevettori $F: \nonstandard{U} \longrightarrow \nonstandard{V}$ \emph{realizza} un campo vettoriale classico $ X $ se per ogni funzione differenziabile \begin{equation*}
        h : U \longrightarrow \R
    \end{equation*}si ha che \begin{equation*}
        X\,h(a) = \vec{a\, F(a)}\,h,\qquad \forall\, a \in U.
    \end{equation*} 
}{kcjndkjndjdjjdjjjjjjjjjjj}{}%Un vettore tangente, come abbiamo visto, associa a ciascuna funzione f : M to R un numero: il valore della corrispondente derivata direzionale di f in quel punto. Un campo vettoriale X,  pertanto, associa alla funzione f una nuova funzione su M, che denotiamo con X(f).

Questa definizione prende in considerazione solo i punti standard.

\paragrafo{Proposizione}{%
    Sia $ p \in M $, $ V $ un intorno coordinato di $ p $. Per ogni campo vettoriale \begin{equation*}
        X: V \longrightarrow \R^{n}
    \end{equation*}esiste un campo locale di prevettori che \emph{realizza} $ X $, definito in un intorno di $ p $.
}{dkjnckjdjdjdjjjjjdjjjdicoij}{}
\dimframmento{dkjnckjdjdjdjjjjjdjjjdicoij}{
    Sia $ V' \subseteq \R^{n}$ l'immagine attraverso le coordinate di $ V $, e sia \begin{equation*}
        X':V' \longrightarrow \R^{n}
    \end{equation*}la rappresentazione in coordinate di $ X $. 
    
    Sia $ U' \subseteq V' $ un intorno dell'immagine di $ p $ tale che: \begin{itemize}
        \item $ \overline{U'} $ è compatto;
        \item $ \overline{U'} \subseteq V'$.
    \end{itemize}Se con $ X' $ si denota anche l'estensione non standard del campo vettoriale, si ha che per ogni $ a \in \nonstandard{U'} $, $ a + \lambda \cdot X'(a) \in V'$. Infatti, $ \nonstandard{\left(\overline{U'}\right)} = \hhalo{\left(\overline{U'}\right)}$ poiché $ \overline{U'} $ è compatto per ipotesi, e dunque per ogni $ b \in \nonstandard{\left(\overline{U'}\right)} $, $ X'(b) $ è finito. 
    
    Quindi, per ogni $ a \in \nonstandard{U} \subseteq \nonstandard{\left(\overline{U'}\right)} $, $ X'(a) $ è finito, e dunque \begin{equation*}
        a + \lambda \cdot X'(a) \in \halo(a) \subseteq \hhalo{\left(\overline{U'}\right)} =\nonstandard{\left(\overline{U'}\right)} \subseteq \nonstandard{V'}.
    \end{equation*}È possibile ora definire la mappa \begin{align*}
    F':\nonstandard{U'} &\longrightarrow \nonstandard{V'}  \\
    a &\longmapsto a + \lambda \cdot X'(a)
    \end{align*}che induce un campo di prevettori locale $ F:\nonstandard{U} \longrightarrow \nonstandard{V} $ che realizza $ X $ in $ U $.\qed
}

La costruzione del campo di prevettori espresso in coordinate da $ F(a) = a + \lambda X(a) $ presentata nella precedente definizione sarà utilizzata più volte. $ F $ prenderà il nome di \emph{campo di prevettori naturale} di un campo vettoriale $ X $.

\proposizione{kjndjcjdkjncjdjckjnskdkjckskjn}{
    Sia $ U \subseteq \R^{n} $ aperto, e sia $ X: U \longrightarrow \R^{n} $ un campo vettoriale classico; sia $ \varphi: U \longrightarrow W \subseteq \R^{n} $ un cambio di coordinate\footnote{Ovvero almeno di classe $ C^{1} $.}; sia $ Y: W \longrightarrow \R^{n} $ il campo vettoriale corrispondente, ovvero tale per cui \begin{equation*}
        Y\left(\varphi(a)\right) = J_{\varphi}(a) \cdot X(a) 
    \end{equation*}dove $ J_{\varphi}(a)  $ è la matrice jacobiana di $ \varphi $ calcolata in $ a $. Siano $ F,G $ le mappe: \begin{equation*}
        F(a) = a + \lambda\,X(a),\quad G(b) = b+ \lambda\,Y(b),\qquad \forall\, a \in \nonstandard{U}, b \in \nonstandard{W}.
    \end{equation*}Allora $ F\equiv \varphi^{-1} \circ G \circ \varphi $, o equivalentemente, \begin{equation*}
        \varphi \circ F(a) - G\circ\varphi(a)\pprec \lambda,\qquad \forall\, a \in \hhalo{U}.
    \end{equation*}
}

    Questa proposizione assicura che il campo locale di prevettori che realizza $ X $ come costruito nella dimostrazione della Proposizione \framref{dkjnckjdjdjdjjjjjdjjjdicoij} è indipendente dalla scelta di coordinate sull'intorno coordinato di $ p $.

\dimostrazioneprop{kjndjcjdkjncjdjckjnskdkjckskjn}{
    Sia $ a \in \hhalo{U} $ fissato. Siano $ \varphi^{i}, Y^{i} $ e $ G^{i} $, rispettivamente, l'$i$-esima coordinata di $ \varphi $, $ Y $ e $ G $. Sia $ D^{i} $ l'$i$-esima riga di $ J_{\varphi}(a) $. Allora si ha \begin{align*}
        \varphi^{i}\circ F(a) - G^{i}\circ \varphi(a) &= \varphi^{i} \left(a+\lambda\,X(a)\right) - \varphi^{i}(a)-\lambda\, Y^{i}\left(\varphi(a)\right)\\  
        &= \left[\varphi^{i} \left(a+\lambda\,X(a)\right) - \varphi^{i}(a)\right] - \lambda\, D^{i} \cdot X(a).
    \end{align*}Per la Proprietà \framref{dckjndjjjjjjdjscacchina}, esiste $ c \in \left[a,a+\lambda\,X(a)\right] $ tale che, se $ D^{i}_c $ è l'$i$-esima riga di $ J_{\varphi}(c)  $, \begin{equation*}
        \varphi^{i} \left(a+\lambda\,X(a)\right) - \varphi^{i}(a) = \lambda\,D^{i}_c \,X(a)
    \end{equation*}Dunque si ha che \begin{equation*}
        \varphi^{i}\circ F(a) - G^{i}\circ \varphi(a) = \lambda\, \left(D^{i}_c-D^{i}\right) X(a) \pprec \lambda.
    \end{equation*}Poichè questo è vero per ogni componente $ i $, si ha \begin{equation*}
        \varphi \circ F(a) - G\circ\varphi(a)\pprec \lambda,\qquad \forall\, a \in \hhalo{U}.\qedd
    \end{equation*}
}
\definizione{
    Si definisce $ I $ come il campo di prevettori identità, su $ \nonstandard{M} $ o su $ \nonstandard{U} $ per ogni $ U \subseteq M $, per ogni $ a $: $ I(a)=a $.
}
\paragrafo{Osservazione}{%
Il campo di prevettori identità $ I $ corrisponde al campo vettoriale classico zero, tramite la procedura illustrata nella dimostrazione della Proposizione \framref{dkjnckjdjdjdjjjjjdjjjdicoij}.
}{kjndkjncskjndjjjdjjdjdjdjjdsjhbcsjhdbcs}{}
\subsection{Condizioni di regolarità}
\definizione{
    Un campo di prevettori su una varietà differenziabile $ M $ è di classe $ D^{1} $ se per ogni $ a,b \in \hhalo{M} $ e $ (a,b) \subseteq P_{a}M $, allora in coordinate \begin{equation*}
        F(a)-a-F(b)+b \prec \lambda\norma{a-b}.
    \end{equation*}
}
\osservazione{
    Questa definizione è motivata dal modo in cui la condizione di Lipschitzianità su un campo vettoriale si trasmette sul campo di prevettori che lo realizza.

    Infatti, un campo vettoriale $ X: U \longrightarrow \R^{n} $ si dice Lipschitziano se esiste $ K \in \R $ tale per cui, per ogni $ a,b \in U $ \begin{equation}
        \norma{X(a)-X(b)}\le K\,\norma{a-b}\label{eq:unino}
    \end{equation}Sia ora $ F $ il campo di prevettori che realizza $ X $, come nella dimostrazione della Proposizione \framref{dkjnckjdjdjdjjjjjdjjjdicoij}, dove $ F(a)-a= \lambda\,X(a) $. La \eqref{eq:unino} diventa: \begin{equation*}
        \norma{\left(F(a)-a\right)-\left(F(b)-b\right)}\le K\lambda\,\norma{a-b}
    \end{equation*}ovvero esattamente $ F $ di classe $ D^{1} $.
}
\paragrafo{Notazione}{%
    Si utilizzerà la notazione di Eulero per le differenze finite. Siano $ V,W $ spazi vettoriali (di campo $ \R $ o $ \starR $); siano $ A \subseteq V $, $ a \in A $ e $ v_1,\dots,v_{k} \in V $ tali che \begin{equation*}
        a+e_1\,v_1+\dots+e_{k}\,v_{k} \in A  
    \end{equation*}per ogni $ (e_1,\dots,e_{n} ) \in \{0,1\}^{k} $.

    Se $ F:A \longrightarrow W $, si definisce la $ k$-esima differenza divisa $ \Delta_{v_1,\dots,v_{k}}^{k} F(a) $ come segue: \begin{equation*}
        \Delta_{v_1,\dots,v_{k}}^{k} F(a) = \sum_{(e_1,\dots,e_{n} ) \in \{0,1\}^{k}} (-1)^{\sum e_{j} } F(a+e_1\,v_1+\dots+e_{k}\,v_{k})
    \end{equation*}
}{}{}
\osservazione{
    Utilizzando la notazione delle differenze finite, la condizione $ D^{1} $ diventa \begin{equation*}
        \Delta_{b-a}^{1}(F-I)(a)\prec \lambda\norma{a-b} 
    \end{equation*}per ogni $ a,b \in \hhalo{M}$, con $ a-b\prec \lambda $. 
    
    Equivalentemente, se $ v=b-a $, la condizione diventa: $ \displaystyle\Delta_{v}^{1}(F-I)(a)\prec \lambda\norma{v}  $.
}
\osservazione{
    Tramite facili passaggi algebrici, è evidente che per ogni $ k\ge 2 $: \begin{equation*}
        \Delta_{v_1,\dots,v_{k}}^{k} I(a) = 0 
    \end{equation*}per ogni $ v_1,\dots,v_{k}  $ e per ogni $ a $. Questo consente di dare la seguente definizione. 
}
\definizione{
    Un campo di prevettori $ F $ su un intorno $ U $ di una varietà differenziabile $ M $ è di classe $ D^{k} $ se, per ogni $ a \in \hhalo{U} $ e per ogni $ v_1,\dots,v_{k} \in \nonstandard{(\R^n)} $, con $ v_{i} \prec \lambda  $, si ha, in coordinate, \begin{equation*}
        \Delta_{v_1,\dots,v_{k} }^{k} F(a) \prec \lambda\,\norma{v_1} \cdots\norma{v_{k} }. 
    \end{equation*}}

Per semplicità, si dà anche la seguente definizione, rindondante, per i campi di prevettori di classe $ D^{2} $. 
\definizione{
    Un campo di prevettori $ F $ su un intorno $ U $ di una varietà differenziabile $ M $ è di classe $ D^{2} $ se, per ogni $ a \in \hhalo{U} $ e per ogni $ v,w \in \nonstandard{(\R^n)} $, con $ v,w \prec \lambda  $, si ha, in coordinate, \begin{equation*}
        \Delta_{v,w}^{2} F(a) = F(a)-F(a+v)-F(a+w)+F(a+v+w) \prec \lambda\,\norma{v}\norma{w }. 
    \end{equation*}
}
\osservazione{
    Anche se non esplicitamente mostrato in questo lavoro, le definizioni dei campi di prevettori di classe $ D^{1} $ e $ D^{2} $ sono indipendenti dalle coordinate. Per una trattazione più approfondita, si rimanda alla Proposizione 3.10 e alla Proposizione 3.14 di \cite{nowik_differential_2015}.
}

Nella proposizione successiva si mostrerà che il campo di prevettori costruito secondo la dimostrazione della Proposizione \framref{dkjnckjdjdjdjjjjjdjjjdicoij}, indotto da un campo vettoriale di classe $ C^k $, è un campo di prevettori di classe $ D^{k} $. 

\proposizione{dcoickjdkjnckjnkskjnksckjndkjnckjnckjdkjnckkjnksjncj}{
    Sia $ W \subseteq \R^{n} $ un aperto, e sia $ X: W \longrightarrow \R^{n} $ un campo vettoriale classico di classe $ C^{k} $. Allora, per ogni $ a \in \hhalo{W} $ e ogni $ v_1,\dots,v_{k} \in \nonstandard{(\R^n)}  $, con $ v_{i} \pprec 1  $, si ha \begin{equation*}
        \Delta_{v_1,\dots,v_{k} }^{k} X(a) \prec \norma{v_1}\norma{v_2} \cdots\norma{v_{k} }. 
    \end{equation*}
}
\dimostrazioneprop{dcoickjdkjnckjnkskjnksckjndkjnckjnckjdkjnckkjnksjncj}{
    Sia $ a \in \hhalo{W} $ fissato, e sia $ U \subseteq W $ un intorno di $ \st(a) $ più piccolo di $ W $, tale che su $ U $ tutte le $ k $-esime derivate parziali di $ X $ siano limitate. Siano $ X^{1},\dots,X^{n}  $ le componenti di $ X $. Sia $ p \in U$, e siano $ v_1,\dots,v_{k} \in \R^{n}  $ tale che \begin{equation*}
        p + s_1v_1+\dots+s_{k}v_{k} \in U  
    \end{equation*}per ogni $ 0\le s_1,\dots,s_{k} \le 1 $. Si definisce quindi, per ogni $  i = 1,\dots,n $ \begin{equation*}
        \varphi^{i}(s_1,\dots,s_{k} ) = X^{i}(p + s_1v_1+\dots+s_{k}v_{k}).
    \end{equation*}Iterando $ k $ volte il teorema del valor medio si ottiene che esiste $ (t_1,\dots,t_{k} ) \in  [0,1]^{k} $ tale per cui \begin{equation*}
        \sum_{(e_1,\dots,e_{k} ) \in \{0,1\}^{k}} (-1)^{\sum e_{j}  } \varphi^{i} (e_1,\dots,e_{k} ) = \frac{\partial\null^{\hspace{0.6pt}k}}{\partial s_1\ \partial s_2 \dots \partial s_{k} } \varphi^{i}(t_1,\dots,t_{k} )\,(-1)^{k}.
    \end{equation*}Dunque si ha che \begin{align*}
        \abs{\Delta_{v_1,\dots,v_{k} }^{k}X^{i}(p) } &= \abs{\sum_{(e_1,\dots,e_{k} ) \in \{0,1\}^{k}} (-1)^{\sum e_{j}  }X^{i}(p + e_1v_1+\dots+e_{k}v_{k})}\\ 
        &= \abs{\sum_{(e_1,\dots,e_{k} ) \in \{0,1\}^{k}} (-1)^{\sum e_{j}  } \varphi^{i} (e_1,\dots,e_{k} )}\\ 
        &= \abs{\frac{\partial\null^{\hspace{0.6pt}k}}{\partial s_1\ \partial s_2 \dots \partial s_{k} } \varphi^{i}(t_1,\dots,t_{k} )} \le C_{i}\norma{v_1}\cdots\norma{v_{k} } 
    \end{align*}dove $ C_{i} \in \R $ è determinato dalla limitatezza della $ k $-esima derivata parziale di $ X^{i} $ in $U$.

    Poiché questo è vero per ogni $ X^{i} $, $ i=1,\dots,n $, si ha che esiste $ K \in \R $ tale per cui \begin{equation*}
        \norma{\Delta_{v_1,\dots,v_{k} }^{k}X(p)}\le K \norma{v_1}\norma{v_2} \cdots\norma{v_{k} }
    \end{equation*}per ogni $ p \in U $ e per ogni $ v_1,\dots,v_{k} \in \R^{n} $ tale per cui $ p + s_1v_1+\dots+s_{k}v_{k} \in U   $ per ogni $ 0\le s_1,\dots,s_{k} \le 1  $.

    Per il principio di transfer questo è vero, con lo stesso $ K $, per ogni $ p \in \nonstandard{U} $ e per ogni $ v_1,\dots,v_{k} \in \nonstandard{(\R^n)}  $ tali per cui $ a+s_1v_1+\dots+s_{k}v_{k} \in \nonstandard{U} $ per ogni $ 0\le s_1,\dots,s_{k} \le 1  $, $ s_{i} \in \starR  $. In particolare, questo è vero per $ a $ e per ogni $ v_1,\dots,v_{k} \in \nonstandard{(\R^n)}$ con $ v_{i} \pprec 1 $. \qed
}
\proposizione{kjndjckjndjckjnskdkjncjdkskjncjdkjckjnjnjnjckjnkskjnc}{
    Sia $ F $ un campo di prevettori di classe $ D^{1} $, e siano $ a, b \in \hhalo{M} $, con $ a-b \prec \lambda $, allora, in coordinate, \begin{equation*}
        \norma{a-b} \prec \norma{F(a)-F(b)}\prec \norma{a-b}.
    \end{equation*}Più in dettaglio, dato $ K \prec 1 $ tale che, in coordinate, $ \norma{F(a)-F(b)-a+b} \le K \lambda \norma{a-b} $, si ha \begin{equation*}
        (1-K\lambda)\,\norma{a-b}\le \norma{F(a)-F(b)}\le (1+K \lambda)\,\norma{a-b}
    \end{equation*}
}
\dimostrazioneprop{kjndjckjndjckjnskdkjncjdkskjncjdkjckjnjnjnjckjnkskjnc}{
    Segue dalla disuguaglianza triangolare: \begin{equation*}
        \big|\norma{F(a)-F(b)}-\norma{a-b}\big| \le \norma{F(a)-F(b)-a+b}
    \end{equation*}e per ipotesi $ \norma{F(a)-F(b)-a+b}\le K \lambda \norma{a-b} $, da cui segue la tesi.\qed
}

Si dimostrano ora un lemma tecnico, che servirà, tra le altre cose, a dimostrare che un campo di prevettori di classe $ D^{2} $ è anche di classe $ D^{1} $.
\paragrafo{Lemma}{%
    Sia $ B \subseteq \R^{n} $ una palla aperta centrata nell'origine, e sia $ a \in \halo(0) $, e sia $ 0\neq v \in \nonstandard{(\R^n)} $ con $ v\pprec 1 $. Se $ G: \nonstandard{B} \longrightarrow \nonstandard{(\R^n)} $ è una funzione \emph{interna} tale per cui \begin{equation*}
        \Delta_{v,v}^{2}G(x) \prec \norma{v}\,\norma{G(a)-G(a+v)} 
    \end{equation*}per ogni $ x \in \hhalo{B} $, allora esiste $ m \in \nonstandard{ \N} $ tale per cui $ a+mv \in \hhalo{B} $ e \begin{equation*}
        G(a)-G(a+v) \prec \norma{v}\,\norma{G(a)-G(a+mv)}.
    \end{equation*}
}{djjjjjjjjjframmentinoframmentucciovieniqui}{}
\dimframmento{djjjjjjjjjframmentinoframmentucciovieniqui}{
    Sia $ 0<r \in \R $ un raggio più piccolo di quello di $ B $, e per $ x \in \starR $ si denoti con $ \floor{x} \in \nonstandard{ \N} $ la parte intera di $ x $. Sia dunque $ N = \floor{r/\norma{v}} $. 

    Dunque ogni $ m\le N $ soddisfa la condizione $ a + mv \in \hhalo{B} $. 
    
    Per semplicità di notazione, siano $ A \coloneqq G(a)-G(a+v) $, $ x_{j} \coloneqq a+jv $, per ogni $ 0\le j \le N $. Si ha allora che \begin{equation*}
        G(x_{j} )-2G(x_{j+1} )+G(x_{j+2} ) = \Delta_{v,v}^{2}G(x_{j} ) = C_{j}\norma{v}\,\norma{A}  
    \end{equation*}con $ C_{j} \prec 1 $.

    Sia $ C $ il massimo dei $ C_0,\dots,C_{N}  $: $ C\prec 1 $, e, per ogni $ 0 \le j\le N $ \begin{equation*}
        G(x_{j} )-2G(x_{j+1} )+G(x_{j+2} ) \le C \norma{v}\,\norma{A}. 
    \end{equation*}

    Per $ k \le N $, si ha \begin{multline*}
        \norma{A-\left(G(x_{k} )-G(x_{k+1} )\right)} = \norma{\left(G(x_{0} )-G(x_{1} )\right)-\left(G(x_{k} )-G(x_{k+1} )\right)}\le\\ 
        \le \sum_{j=0}^{k-1} \norma{\left(G(x_{j} )-G(x_{j+1} )\right)-\left(G(x_{j+1} )-G(x_{j+2} )\right)} \le Ck\norma{v}\norma{A}
    \end{multline*}e quindi, per ogni $ m\le N $ \begin{equation*}
        \norma{mA - \left(G(x_0)-G(x_{m} )\right)} = \norma{\sum_{k=0}^{m-1}\left(A-\left(G(x_{k} ) - G(x_{k+1} ) \right)\right) } \le Cm^{2}\norma{v}\norma{A}.
    \end{equation*}Si può quindi dire che \begin{equation*}
        \norma{mA - \left(G(x_0)-G(x_{m} )\right)} =  Km^{2}\norma{v}\norma{A}.
    \end{equation*}per qualche $ K\prec 1 $, e, per la disuguaglianza triangolare, \begin{equation*}
        m\norma{A} - Km^{2}\norma{v}\norma{A}\le \norma{G(x_0)-G(x_{m} )}
    \end{equation*}e, moltiplicando per $ \norma{v} $, $ m\norma{v}\norma{A}(1-Km\norma{v})\le \norma{v}\norma{G(x_0)-G(x_{m} )} $.

    Sia ora $ m=\min\left\{N, \floor{1/(2K\norma{v})}\right\}\le N $; allora $ (1-Km\norma{v})\ge 1/2 $, e dunque \begin{equation*}
        m\norma{v}\norma{A}/2\le \norma{v}\norma{G(x_0)-G(x_{m} )}.
    \end{equation*}Per definizione di $ N $ e poiché $ K $ è finito, $ m\norma{v}/2 $ non è infinitesimo, e dunque \[A\prec \norma{v}\norma{G(x_0)-G(x_{m} )} \qedd\]
}
\proposizione{dkjncjdkkjncsjdkjncsiujdoijfdosoidocodidislkdmlkmlkm}{
    Se $ F $ è un campo di prevettori di classe $ D^{2} $ su $ W $, allora $ F $ è di classe $ D^{1} $.
}
\dimostrazioneprop{dkjncjdkkjncsjdkjncsiujdoijfdosoidocodidislkdmlkmlkm}{
    Sia $ a \in \hhalo{W} $, e sia $ B \subseteq \R^{n} $ una palla aperta centrata, nelle date coordinate, in $ \st({a}) $. Per mostrare che $ F $ è di classe $ D^{1} $, si deve mostrare che, in coordinate, per ogni $ b $ tale che $ b-a \prec \lambda $ vale $ F(a)-a-F(b)+b \prec \lambda\norma{a-b} $. 

    Se $ v \coloneqq b-a $ e $ G=F-I $ (i.e. $ G(x)=F(x)-x $) la condizione è equivalente a richiedere che \begin{equation*}
        G(a)-G(a+v)\prec \lambda \norma{v}.
    \end{equation*}I casi sono due\footnote{Vedi Osservazione \framref{djknckjnsjdjdjcskjnskjnckjnskjnd}.} \begin{itemize}
        \item Se $ G(a)-G(a+v)\pprec \lambda\norma{v} $ la tesi è banalmente verificata.
        \item Altrimenti si ha $ \lambda v \prec \norma{G(a)-G(a+v)} $, e dunque \[\lambda\norma{v}^{2}\prec \norma{v}\norma{G(a)-G(a+v)};\]inoltre, per ipotesi, poiché $ F $ è di classe $ D^{1}$: per ogni $ x $: \begin{equation*}
            \Delta_{v,v}^{2}G(x) = \Delta_{v,v}^{2}  F(x) \prec \lambda\norma{v}^{2}.
        \end{equation*}Quindi $ \Delta_{v,v}^{2}G(x)\prec \norma{G(a)-G(a+v)} $, e per il Lemma \framref{djjjjjjjjjframmentinoframmentucciovieniqui} esiste $ m \in \nonstandard{ \N}$ tale che $ a+mv \in \hhalo{B} $ e {\begin{align*}
            G(a)-G(a+v)&\prec \norma{v}\norma{G(a)-G(a+mv)} \\&\le \norma{v}\left(\norma{G(a)}+\norma{G(a+mv)}\right)\prec \norma{v}\lambda
        \end{align*}}poiché $ F $ è un campo di prevettori, e dunque $ G(x)=F(x)-x\prec \lambda $ per ogni $ x $.\qed
    \end{itemize}
}
\paragrafo{Osservazione}{%
Si enunciano di seguito due risultati, che non saranno dimostrati, in merito alla composizione di campi di prevettori. Si veda la Proposizione 3.16 e la Proposizione 3.17 di \cite{nowik_differential_2015}.\begin{itemize}
    \item Se $ F,G $ sono campi di prevettori di classe $ D^{1} $, allora $ F\circ G $ è un campo di prevettori di classe $ D^{1} $. 
    \item Se $ F,G $ sono campi di prevettori di classe $ D^{2} $, allora $ F\circ G $ è un campo di prevettori di classe $ D^{2} $.
\end{itemize}
}{dkcjnacjdjdkjcnkjnskjncjdjckjnsjcjdkjnc}{}

\section{Proprietà globali dei campi di prevettori}

\proposizione{lkmdkclslkmdkcoijskiklkmclknsdlkmcslkmslkmlkm}{
    Sia $ F $ un campo di prevettori.\begin{enumerate}
        \item Se $ W $ è un intorno coordinato con immagine $ U \subseteq \R^{n} $, e $ B \subseteq U $ è una palla chiusa, allora esiste un $ C $ \emph{finito} tale per cui \begin{equation*}
            \norma{F(a)-a}\le C \lambda
        \end{equation*}per ogni $ a \in \nonstandard{B} $. (Si noti che $ F $ denota sia il campo di prevettori che la sua azione in coordinate).
        \item Se $ G $ è un altro campo di prevettori, allora esiste un $ \beta $ finito tale per cui \begin{equation*}
            \norma{F(a)-G(a)}\le \beta\lambda
        \end{equation*}per ogni $ a \in \nonstandard{B} $.
        \item Se inoltre $ F\equiv G $, allora esiste un tale $ \beta $ \emph{infinitesimale}.
    \end{enumerate}
}
\dimostrazioneprop{lkmdkclslkmdkcoijskiklkmclknsdlkmcslkmslkmlkm}{
    \begin{enumerate}
        \item È conseguenza di 2, considerando $ G(a)=a $ per ogni $ a $. 
        \item Sia $\displaystyle 
            A=\left\{n \in \nonstandard{ \N}: \norma{F(a)-G(a)}\le n\lambda\text{ per ogni }a \in \nonstandard{B}\right\}$. 

        Ogni $ n \in \nonstandard{ \N} $ infinito è contenuto in $ A $, e pertanto, per il principio dell'\straniero{underspill} esiste $ C \in A $ \emph{finito}. Questo porta alla tesi.
        \item Sia ora invece \begin{equation*}
            A=\left\{n \in \nonstandard{ \N}: \norma{F(a)-G(a)}\le \frac{\lambda}{n}\lambda\text{ per ogni }a \in \nonstandard{B}\right\}.
        \end{equation*}Ogni $ n \in \nonstandard{ \N} $ finito è contenuto in $ A $, e perciò, per \straniero{overspill}\footnote{Posso applicarlo in quanto $ A $ è un insieme interno per il Teorema \teoref{jdjdjjjjjjdjjjkjn}.} esiste un $ m \in \nonstandard{ \N} $ infinito contenuto in $ A $. Pertanto, sia $ \beta = 1/m \pprec 1 $.\qed
    \end{enumerate}
}
\proposizione{dlikjckjdlkmnckslkmdkcoslkdlkj}{
    Sia $ F $ un campo di prevettori di classe $ D^{1} $. Se $ W $ è un intorno coordinato con immagine $ U \subseteq \R^{n} $ e $ B \subseteq U $ è una palla chiusa, allora esiste $ K \in \R $ tale per cui \begin{equation}
        \norma{F(a)-a-F(b)+b}\le K\lambda\norma{a-b}.\label{eq:chissa}
    \end{equation}per ogni $ a,b \in \nonstandard{B} $.
    
    Segue che \begin{equation}
        (1-K\lambda)\norma{a-b}\le \norma{F(a)-F(b)}\le (1+K\lambda)\norma{a-b}.\label{eq:minbound}
    \end{equation}
}
\dimostrazioneprop{dlikjckjdlkmnckslkmdkcoslkdlkj}{
    Sia $ N=\floor{1/\lambda} $. Dati $ a,b \in \nonstandard{B} $, per $ k=0,\dots,N $, sia \begin{equation*}
        a_{k} = a + \frac{k}{N} (b-a),
    \end{equation*}allora $ a_{k}-a_{k+1} \prec \lambda  $. 

    Sia dunque $ C_{ab}  $ il massimo della seguente espressione, al variare di $ k=0,\dots,N-1 $: \begin{equation*}
        \frac{\norma{F(a_{k} )-a_{k}-F(a_{k+1} )+a_{k+1}  }}{\lambda \norma{a_{k}-a_{k+1}  }}
    \end{equation*}che è finito, poiché $ F $ è di classe $ D^{1} $. Dunque, per ogni $ 0\le k \le N-1 $, si ha \begin{equation*}
        \norma{F(a_{k} )-a_{k}-F(a_{k+1} )+a_{k}  }\le C_{ab}\lambda\norma{a_{k}-a_{k+1}  } = C_{ab}\lambda\frac{\norma{a-b}}{N}.  
    \end{equation*}Segue che \begin{equation*}
        \norma{F(a)-a-F(b)+b}\le \sum_{k=0}^{N-1}\norma{F(a_{k})-a_{k}-F(a_{k+1}  )+a_{k+1}  } \le C_{ab}\lambda\norma{a-b}. 
    \end{equation*}Se ora si definisce $ A $ come \begin{equation*}
        A= \set{n \in \nonstandard{ \N}:  \norma{F(a)-a-F(b)+b}\le  n\lambda\norma{a-b}\text{ per ogni }a,b \in \nonstandard{B} },
    \end{equation*}poiché tutti gli $ C_{ab} $ sono finiti, si ha che $ A $ contiene ogni $ n \in \nonstandard{ \N} $ infinito, e pertanto, per \straniero{underspill}, esiste $ K \in A $ finito, e quindi \eqref{eq:chissa} vale.\qed
}
\proposizione{dkjncjdkjncjdkjncjdkjncjdjcjdkjncsdkjnkj}{
    Sia $ F $ un campo di prevettori di classe $ D^{2} $. Se $ W $ è un intorno coordinato con immagine $ U \subseteq \R^{n} $ e $ B \subseteq U $ è una palla chiusa, allora esiste $ K \in \R $ tale per cui \begin{equation*}
        \norma{\Delta_{v,w}^{2} F(a) } \le K\lambda\norma{w}\norma{v}
    \end{equation*}per ogni $ a \in \nonstandard{B} $ e per ogni $ v,w \in \nonstandard{(\R^n)}$ tali per cui \begin{equation*}
        a+v,a+w,a+v+w \in \nonstandard{B}.
    \end{equation*}
}
\dimostrazioneprop{dkjncjdkjncjdkjncjdkjncjdjcjdkjncsdkjnkj}{
    Sia $ N = \floor{1/\lambda} $. Dato $ a \in \nonstandard{B} $ e $ v,w \in \starR$ tali per cui valgano le ipotesi, sia \begin{equation*}
        a_{k,l} = a+ \frac{k}{N} v + \frac{l}{N} w 
    \end{equation*}per $ 0 \le k,l\le N $. Sia $ C_{avw} $ il massimo della seguente espressione al variare di $ k,l $, $ 0\le k,l\le N-1 $: \begin{equation*}
        \frac{
            \norma{\Delta_{v/N, w/N }^{2} F(a_{k,l} ) }
        }{\lambda\norma{v/N}\norma{w/N}}.
    \end{equation*}Poiché $ F $ è di classe $ D^{2} $, $ C_{avw}$ è finito. Dunque, per ogni $ 0\le k,l\le N-1 $: \begin{align*}
        \norma{F(a_{k,l} )-F(a_{k+1,l} ) - F(a_{k,l+1} ) + F(a_{k+1,l+1} ) }&= \\ = \norma{\Delta_{v/N, w/N }^{2} F(a_{k,l} )} &\le C_{avw} \lambda\norma{v/N}\norma{w/N}
    \end{align*}Con un argomentazione simile a quelle delle dimostrazioni precedenti, si ha che \begin{equation*}
        \norma{F(a)-F(a+v)-F(a+w)+F(a+v+w)} \le C_{avw}\lambda\norma{v}\norma{w},
    \end{equation*}e, utilizzando il principio dell'\straniero{underspill}, esiste un $ K $ finito tale per cui vale la tesi.\qed
}
\proposizione{djknckjndjckjnxkjnkjnckjndskjnckjnkjn}{
    Se $ F $ è un campo di prevettori di classe $ D^{1} $, allora è iniettivo su $ \hhalo{M} $.
}
\dimostrazioneprop{djknckjndjckjnxkjnkjnckjndskjnckjnkjn}{
    Siano $ a\neq b \in \hhalo{M} $.\begin{itemize}
        \item Se $ \st(a)\neq \st(b) $, allora chiaramente\footnote{Segue dalla \eqref{eq:minbound} con $ \norma{a-b} $ non infinitesimo.}si avrà $ F(a) \neq F(b) $;
        \item se invece $ \st(a) = \st(b) $, allora esiste una palla chiusa che soddisfi le ipotesi della Proposizione \framref{dlikjckjdlkmnckslkmdkcoslkdlkj}; pertanto, vale l'equazione \eqref{eq:minbound}.\qed
    \end{itemize}
}
\proposizione{djncsjkndkjnckjsndkjncskjndkjncskjndjjjdj}{
    Siano $ B_1 \subseteq B_2 \subseteq B_3 \subseteq \R^{n} $ delle palle chiuse centrate nell'origine, di raggi $ r_1<r_2<r_3 $. Se $ F: \nonstandard{B_2} \longrightarrow \nonstandard{B_3} $ è un campo di prevettori locale di classe $ D^{1} $, allora $ F(\nonstandard{B_2}) \supseteq \nonstandard{B_1} $.
}
\dimostrazioneprop{djncsjkndkjnckjsndkjncskjndkjncskjndjjjdj}{
    Sia $ 0<s \in \R $ fissato, $ s<r_2-r_1 $ e $ s<r_3-r_2 $. Si applicherà il principio di transfer al seguente enunciato: per ogni funzione $ f: B_2 \longrightarrow B_3 $, se $ \norma{f(x)-f(y)}\le 2\norma{x-y} $ per ogni $ x,y \in B_2 $ e $ \norma{f(x)-x}<s $ per ogni $ x \in B_2 $, allora $ f(B_2) \supseteq B_1 $.

    L'enunciato è vero poiché la condizione su $ f $ implica che questa sia continua. La continuità, unita alla definizione di $ s $, garantisce che per ogni $ x \in \partial B_2 $ il segmento $ [x,f(x)] $ sia contenuto in $ B_3\setminus B_1 $. 

    Segue che $ \restriction{f}{\partial B_2}: \partial B_2  \longrightarrow B_3\setminus B_1$ sia omotopa alla inclusione di $ \partial B_2 $ all'interno di $ B_3\setminus B_1 $, con omotopia \begin{align*}
    H:\partial B_2\times [0,1] &\longrightarrow B_3\setminus B_1 \\
    (x,t) &\longmapsto tx+(1-t)f(x).
    \end{align*}Dunque, se per assurdo esistesse $ p \in B_1 $ non compreso in $ f(B_2) $, allora $ \restriction{f}{\partial B_2} $ sarebbe omotopicamente equivalente alla funzione costante, e pertanto questo sarebbe vero anche per l'inclusione di $ \partial B_2 $. Questo genera contraddizione. %ATOD capire il perché

    Applicando il transfer si ottiene dunque che per ogni funzione interna $ f: \nonstandard{B_2} \longrightarrow \nonstandard{B_3} $, se \begin{align*}
        \norma{f(x)-f(y)}&\le  2\norma{x-y} & \text{per ogni }x,y &\in \nonstandard{B_2},\\ 
        \norma{f(x)-x}&<s & \text{per ogni } x & \in \nonstandard{B_2},
    \end{align*}allora $ f(\nonstandard{B_2})\supseteq \nonstandard{B_1} $. Per la Proposizione \framref{dlikjckjdlkmnckslkmdkcoslkdlkj} un campo di prevettori $ D^{1} $ rispetta le ipotesi; pertanto la proposizione è dimostrata.\qed %ATOD capire
}
\teorema{lddlkmckdkdkclkmckskdlkmclslkmdlkckskdlkmlkcmlkmlslkmdlkmcs}{
   Se $ F: \nonstandard{M} \longrightarrow \nonstandard{M} $ è un campo di prevettori di classe $ D^{1} $, allora \begin{equation*}
       \restriction{F}{\hhalo{M}}:\hhalo{M} \longrightarrow \hhalo{M}.
   \end{equation*}è biiettiva.
}

Il teorema segue dalla Proposizione \framref{djknckjndjckjnxkjnkjnckjndskjnckjnkjn} e dalla Proposizione \framref{djncsjkndkjnckjsndkjncskjndkjncskjndjjjdj}.
\paragrafo{Osservazione}{%
   Se $ M $ non è compatta, allora per un campo di prevettori $ F $ di classe $ D^{1} $ la mappa $ \restriction{f}{\hhalo{M}}:\hhalo{M} \longrightarrow \hhalo{M} $ e la sua inversa $ (\restriction{F}{\hhalo{M}})^{-1} $ non sono interne, poiché il loro dominio non è interno. 

   D'altra parte, però, per ogni $ A \subseteq M $, $ \restriction{F}{\nonstandard{A}} $ è interna. Inoltre, su $ \nonstandard{B_1} $ della Proposizione \framref{djncsjkndkjnckjsndkjncskjndkjncskjndjjjdj}, $ F $ ha una inversa $ F^{-1}:\nonstandard{B_1} \longrightarrow \nonstandard{B_2} $ (ovvero $ F \circ F^{-1} (a) = a $ per ogni $ a \in \nonstandard{B_1} $) e $ F^{-1} $ è interna. 

   Dunque, a meno di ridimensionare il dominio, è sempre possibile assumere che un campo di prevettori locale $ F: \nonstandard{U} \longrightarrow \nonstandard{V} $ di classe $ D^{1} $ ammetta un inverso $ F^{-1}:\nonstandard{U} \longrightarrow \nonstandard{V} $.
}{remarkfourpointeight}{}
\proposizione{kjndjckjnsdkjncjdjckjnsjdkcnskjncjdjdjcjskjnckskjnd}{
    Se $ F $ è di classe $ D^{1} $, allora $ F^{-1} $ è di classe $ D^{1} $. 

    Nello specifico, se per $ x = F^{-1}(a) $, $ y= F^{-1}(b) $ esiste un $ K\prec 1 $ tale per cui \begin{equation*}
        \norma{F(x)-F(y)-x+y}\le K \lambda\norma{x-y}
    \end{equation*}allora, per $ K'=K/(1-K\lambda) $, \begin{equation*}
        \norma{F^{-1}(a)-F^{-1}(b)-a+b}\le K'\lambda\norma{a-b}.
    \end{equation*}
}
\dimostrazioneprop{kjndjckjnsdkjncjdjckjnsjdkcnskjncjdjdjcjskjnckskjnd}{
    Siano $x = F^{-1}(a) $, $ y = F^{-1}(b) $. Allora \begin{gather*}
        \norma{F^{-1}(a)-F^{-1}(b)-a+b} = \\
        = \norma{x-y-F(x)+F(y)} \le K \lambda\norma{x-y} \le K' \lambda\norma{F(x)-F(y)} = K'\lambda\norma{a-b}
    \end{gather*}per la Proposizione \framref{kjndjckjndjckjnskdkjncjdkskjncjdkjckjnjnjnjckjnkskjnc}.\qed
}
\lemma{kjndkjncjdkjnckskdkjnjcsnkjn}{
    Siano $ F,G $ dei campi di prevettori di classe $ D^{1} $. Se $ F(a)\equiv G(a) $ per ogni $ a $ standard, allora $ F(b)\equiv G(b) $ per ogni $ b $ quasi-standard: ovvero $ F\equiv G $.
}
\dimostrazionelem{kjndkjncjdkjnckskdkjnjcsnkjn}{
    Fissato $ b $, sia $ K $ dato dalla Proposizione \framref{dlikjckjdlkmnckslkmdkcoslkdlkj} applicata sia ad $ F $ che a $ G $ in una palla di centro $ a=\st(b) $. Allora \begin{align*}
        \norma{F(b)-G(b)} &= \norma{F(b)-F(a)-b+a+F(a)-G(a)+G(a)-G(b)-a+b}\\ 
        &\le \norma{F(b)-F(a)-b+a} + \norma{F(a)-G(a)} + \norma{G(a)-G(b)-a+b}\\ 
        &\le K \lambda \norma{a-b}+\norma{F(a)-G(a)} + K\lambda\norma{a-b} \pprec \lambda.\qedd
    \end{align*}
}

    Questo lemma garantisce che la definizione di un campo di prevettori utilizzando unicamente i punti standard determina univocamente il campo di prevettori di classe $ D^1 $ a meno di equivalenze.

\section{Flusso di un campo di prevettori}
\definizione{%
Si considerano due casi.
\begin{itemize}
    \item Sia $ F: \nonstandard{M} \longrightarrow \nonstandard{M} $ un campo di prevettori globale, e, per $ 0\le t \in \nonstandard{R} $, sia $ n = n(t) = \floor{t/\lambda} $. Per ogni $ a \in \nonstandard{M} $ si definisce il \emph{flusso di $ F $ al tempo $ t $}, $ F_{t}(a) = F^{n}(a) $. 

    Per definire $ F^{n}(a) $, sia $ \Phi $ la mappa\footnote{Si ricorda che $ \operatorname{Map}(M) $ è l'insieme di tutte le mappe $ f: M \longrightarrow M $} \begin{align*}
    \Phi:\operatorname{Map}(M) \times \N &\longrightarrow \operatorname{Map}(M) \\
    (f,n) &\longmapsto \underset{n\text{ volte}}{\underbrace{f\circ \dots\circ f}}
    \end{align*}e se ne consideri l'estensione non standard $ \nonstandard{\Phi}: \nonstandard{\operatorname{Map}(M)}\times \nonstandard{\N} \longrightarrow \nonstandard{\operatorname{Map}(M)} $: $ F^{n} = \nonstandard{\Phi}(F,n) $.
    \item Sia $ F:\nonstandard{U} \longrightarrow \nonstandard{V} $ un campo di prevettori locale. Si estenda $ F $ ad $ F': \nonstandard{V} \longrightarrow \nonstandard{V} $ definendo $ F'(a) = a $ per ogni $ a \in \nonstandard{V} \setminus \nonstandard{U}$. La mappa $ F' $ può ora essere considerato un campo di prevettori globale su $ \nonstandard{V} $.
    
    Sia $ Y_{n} \coloneqq \set{a \in \nonstandard{U}: (F')^{n}(a) \in \nonstandard{U}}$. Si definisce dunque il \emph{flusso di $ F $ al tempo $ t $} la funzione \begin{align*}
    F_{t} : Y_{n(t)}  &\longrightarrow \nonstandard{U} \\
    a &\longmapsto F_{t}'(a). 
    \end{align*}
\end{itemize}
}
\osservazione{
    È possibile estendere la definizione di $ F_{t} $ (di cui sopra) per $ t\le 0 $. In particolare: \begin{itemize}
        \item per campi di prevettori globali che sono biiettivi su $ \nonstandard{M} $,
        \item per campi di prevettori locali biiettivi nel senso dell'Osservazione \framref{remarkfourpointeight}
    \end{itemize}
    si definisce $ F_{t}  $ per $ t\le  0 $ come $ \left(F^{-1}\right)_{-t} $.
}

\subsection{Dipendenza dalle condizioni iniziali e dal campo di prevettori}
\teorema{dlkmncldlslkmckdlclkmsdkckdkdkkkkkdkk}{
    Sia $ F $ un campo di prevettori locale su $ \nonstandard{U} $ di classe $ D^{1} $. Dato $ p \in U $ e un intorno coordinato di $ p $ di immagine $ W \subseteq \R^{n} $, siano $ B' \subseteq B \subseteq W$ palle chiuse di raggi $ r/2 $ ed $ r $ centrate nell'immagine di $ p $.

    Sia $ K $ una costante finita\footnote{Tale costante esite per la Proposizione \framref{dlikjckjdlkmnckslkmdkcoslkdlkj}.} tale per cui \begin{equation*}
        \norma{F(a)-a-F(b)+b}\le K \lambda\norma{a-b}
    \end{equation*}per ogni $ a,b \in \nonstandard{B}$: allora esiste $ 0<T \in R $ tale per cui $ F_{t}(a) \in \nonstandard{B}$ per ogni $ a \in \nonstandard{B'} $ e $ -T\le t\le T $.

    Inoltre, per ogni $ a,b \in \nonstandard{B'} $ e ogni $ 0\le t\le T $: $ \norma{F_{t}(a) -F_{t}(b) }\le e^{Kt}\norma{a-b} $.

    Prendendo $ K' = K/(1-K\lambda)^{2} $, allora per ogni $ a,b \in \nonstandard{B'}$ e $ -T\le t\le T $: \begin{equation*}
        e^{-K'\abs{t}}\norma{a-b}\le \norma{F_{t}(a)-F_{t}(b)} \le e^{K'\abs{t}}\norma{a-b}.
    \end{equation*}
}
\dimostrazione{dlkmncldlslkmckdlclkmsdkckdkdkkkkkdkk}{
    Si dimostra inizialmente il teorema per $ t\ge 0 $. Sia $ C $ come per la Proposizione \framref{lkmdkclslkmdkcoijskiklkmclknsdlkmcslkmslkmlkm}(1). Per induzione interna\footnote{Si ha che se $ A $ è un sottoinsieme interno di $ \nonstandard{\N} $ chiuso rispetto alla funzione successore $ n \mapsto n+1$, allora $ \nonstandard{N} = A $.} su $ n $ si dimostra che per ogni $ n \in \nonstandard{\N} $: \begin{itemize}
        \item per ogni $ a \in \nonstandard{B'} $ \begin{equation}
            \norma{F^{n}(a)-a}\le \sum_{m=1}^{n} \norma{F^{m}(a)-F^{m-1}(a)} \le nC\lambda\le \frac{r}{2}\label{eq:cosa1}
        \end{equation}e dunque $ F^{n}(a) \in \nonstandard{B} $;
        \item per ogni $ a,b \in \nonstandard{B} $, per la Proposizione \framref{dlikjckjdlkmnckslkmdkcoslkdlkj}, si ha che \begin{equation*}
            (1-K\lambda)\norma{a-b} \le \norma{F(a)-F(b)}\le (1+K\lambda)\norma{a-b},
        \end{equation*}e quindi \begin{equation}
            (1-K\lambda)^{n}\norma{a-b}\le \norma{F^{n}(a)-F^{n}(b)}\le(1+K\lambda)^{n}\norma{a-b}.\label{eq:cosa2}
        \end{equation}
    \end{itemize}
    
    Sia dunque $ T=\frac{r}{2C} $, allora $ 0<T \in \R $. In particolare, \eqref{eq:cosa1} e \eqref{eq:cosa2} valgono per $ n=\floor{t/\lambda} $ per $ 0\le t\le T $, ed è quindi possibile scrivere \begin{align*}
        F_{t}(a) \in \nonstandard{B}\qquad&\text{per ogni }a \in \nonstandard{B'};\\
        \norma{F_{t}(a)-F_{t}(b)} \le (1+K\lambda)^{n}\norma{a-b}\qquad&\text{per ogni }a, b \in \nonstandard{B}.
    \end{align*}

    Segue da banali passaggi algebrici\footnote{%
        Infatti si ha che $ 1+K\lambda \le e^{K \lambda} $, e pertanto $ (1+K\lambda)^{n}\le e^{Kt} $.
    } che per ogni $ a, b \in \nonstandard{B'} \subseteq \nonstandard{B} $ e ogni  $ 0\le t\le T $\begin{equation*}
        \norma{F_{t}(a) -F_{t}(b) }\le e^{Kt}\norma{a-b}\le e^{K'\abs{t}}\norma{a-b}.
    \end{equation*}in quanto se $ K'>K $ allora la $ e^{K'\abs{t}}> e^{K\abs{t}}$. 

    Inoltre, se $ h \coloneqq \frac{1}{1-K\lambda} $, allora\footnote{Scrivendo l'esponenziale come serie di Taylor, risulta chiaro che \begin{equation*}
        e^{-hK\lambda}\le 1-hK\lambda+\frac{h^{2}K^{2}\lambda^{2}}{2} \le 1-hK\lambda+hK^{2}\lambda^{2} = 1-K\lambda.
    \end{equation*}} $ e^{-hKt} \le (1-K\lambda)^{n}$, e pertanto, siccome $ K' = h^{2}K > hK $ si ha che $ e^{-K'\abs{t}} \le e^{-hKt}$ per ogni $ 0\le t\le T $, e dunque \begin{equation*}
        e^{-K'\abs{t}}\norma{a-b}\le \norma{F_{t}(a)-F_{t}(b)  }
    \end{equation*}per ogni $ a,b \in \nonstandard{B'} $. Si è dimostrato il teorema per $ t\ge 0 $.

    Si dimostra ora il teorema per $ t\le 0 $. Si ricorda che per tempi negativi, $ F_{t} \coloneqq (F^{-1})_{\abs{t}} $. 

    A meno di modificare la costante $ C $ precedente (rendendola maggiore) si può supporre che \begin{equation*}
        \norma{F^{-1}(a)-a} \le C\lambda
    \end{equation*}Per la Proposizione \framref{kjndjckjnsdkjncjdjckjnsjdkcnskjncjdjdjcjskjnckskjnd} lo stesso argomento usato per $ F $ si può usare per $ F^{-1} $, a patto di sostituire $ K $ con $ hK $. Pertanto il teorema è dimostrato.\qed
}
\teorema{dlkmckdkclkmskdkclkmskdkcdkkkkkkkk}{
    Sia $ F $ un campo di prevettori locale su $ \nonstandard{U} $, di classe $ D^{1} $, e sia $ G $ un qualsiasi campo di prevettori su $ \nonstandard{U} $. Dato un intorno coordinato compreso in $ U $ con immagine $ W \subseteq \R^{n} $, siano $ A' \subseteq A \subseteq \nonstandard{W} $ insiemi interni. Si supponga: \begin{enumerate}
        \item $ \norma{F(a)-G(a)}\le \beta\lambda $ per ogni $ a \in A $, per qualche costante $ \beta $. (Se $ F\equiv G $ allora $ \beta $ infinitesimale esiste, come assicurato dalla Proposizione \framref{lkmdkclslkmdkcoijskiklkmclknsdlkmcslkmslkmlkm}).
        \item $ \norma{F(a)-a-F(b)+b}\le K\lambda\norma{a-b} $ per ogni $ a,b \in A $, con $ K $ costante finita. (Tale costante esiste per la Proposizione \framref{dlikjckjdlkmnckslkmdkcoslkdlkj}).
        \item $ 0<T \in \R $ è tale per cui $ F_{t}(a)  $ e $ G_{t}(a)$ sono contenuti in $ A $ per ogni $ a \in A' $ e ogni $ 0\le t< T $. 
    \end{enumerate}

    Allora per ogni $ a \in A' $ e $ 0\le t\le T $, \begin{equation*}
        \norma{F_{t}(a)-G_{t}(a)  }\le \frac{\beta}{K}\left(e^{Kt}-1\right)\le \beta t e^{Kt}.
    \end{equation*}
}

\paragrafo{Corollario}{%
    Se $ F \equiv G $, allora $ F_{t}(a)\approx G_{t}(a) $ per ogni $ 0\le t\le T $.
}{djdkjncjdjndimostrazioncinainainainaianinaiai}{}
\dimostrazione{dlkmckdkclkmskdkclkmskdkcdkkkkkkkk}{
    Si dimostra per induzione interna che \begin{equation*}
        \norma{F^{n}(a)-G^{n}(a)}\le \frac{\beta}{K}\left((1+K\lambda)^{n}-1\right)
    \end{equation*}per ogni $ n \in\nonstandard{\N} $. Questo, per l'argomento utilizzato nella dimostrazione del \teoref{dlkmncldlslkmckdlclkmsdkckdkdkkkkkdkk}, dimostra la tesi. Il passo base è vero per ipotesi, e il passo induttivo è dimostrato dalla seguente disuguaglianza: \begin{align*}
        \norma{F^{n+1}(a)-G^{n+1}(a)}&\le \norma{F\left(F^{n}(a)\right)-F\left(G^{n}(a)\right)} + \norma{F\left(G^{n}(a)\right)-G\left(G^{n}(a)\right)}\\ 
        &\le (1+K\lambda)\norma{F^{n}(a)-G^{n}(a)} + \beta \lambda
    \end{align*}dove l'ultimo passaggio è giustificato dalla Proposizione \framref{dlikjckjdlkmnckslkmdkcoslkdlkj}. \qed
}
\definizione{
    Sia $ F $ un campo di prevettori locale di classe $ D^{1} $. In un intorno $ B' \subseteq \R^{n}$ e in un intervallo $ [-T,T] $, come per il Teorema \teoref{dlkmckdkclkmskdkclkmskdkcdkkkkkkkk}, si definisce il \emph{flusso standard indotto da $ F $}: \begin{align*}
    h_{t}^{F}:B' &\longrightarrow M \\
    x &\longmapsto \st\left(F_{t}(x) \right).
    \end{align*}
}
\teorema{kjndjcjdkjncjskjncjsdjcjdjkjnckjndjncjsjdjj}{
    Sia $ F $ un campo di prevettori di classe $ D^{1}$. Allora:\begin{enumerate}
        \item $ h_{t}^{F} $ è continua e Lipschitziana, con costante $ e^{K\abs{t}} $;
        \item $ h_{t}^{F} $ è iniettiva;
        \item se $ G $ è un altro campo di prevettori di classe $ D^{1} $ e $ F\equiv G $, allora $ h_{t}^{F} = h_{t}^{G}$.
    \end{enumerate}
}
\dimostrazione{kjndjcjdkjncjskjncjsdjcjdjkjnckjndjncjsjdjj}{
    Il teorema segue banalmente dal Teorema \teoref{dlkmckdkclkmskdkclkmskdkcdkkkkkkkk} e dal Corollario \framref{djdkjncjdjndimostrazioncinainainainaianinaiai}.\qed
}
\paragrafo{Osservazione}{%
    È possibile dimostrare che se $ F $ è un campo di prevettori ottenuto a partire da un campo vettoriale classico $ X $ tramite la procedura illustrata nella dimostrazione della Proposizione \framref{dkjnckjdjdjdjjjjjdjjjdicoij}, allora $ h_{t}^{F} $ è esattamente il flusso di $ X $ nel senso classico. Si rimanda all'Osservazione 5.8 di \cite{nowik_differential_2015} e al Teorema 14.1 di \cite{keisler_foundation_2007}.
}{osservsADJKNCkjnadkjhdbcakjfhakjdhfgaskjdfghadsfjgh}{}
\subsection{Il campo di prevettori rappresentativo canonico}

\definizione{
    Sia $ F $ un campo di prevettori. Si definisce il campo di prevettori \begin{equation*}
        \tilde{F} = h_{\lambda}^{F} 
    \end{equation*}l'estensione ad un dominio non standard della funzione $ h_{t}^{F}$.
}
\osservazione{
    $\tilde{F}$ è effettivamente un campo di prevettori, ovvero \begin{equation*}
        \tilde{F}(a)-a\prec \lambda
    \end{equation*}per ogni $ a $. Infatti, preso $ C $ come per la Proposizione \framref{lkmdkclslkmdkcoijskiklkmclknsdlkmcslkmslkmlkm}(1), si ha che \begin{equation*}
        \norma{F^{n}(a)-a} \le n C \lambda
    \end{equation*}da cui segue che $ \norma{h^{F}_t(a)-a} \le Ct $. Per il principio di transfer \begin{equation*}
        \norma{\tilde{F}(a)-a} \le C\lambda
    \end{equation*}
}
\lemma{djnckskdkjncjdkjnckjnskjndkjnckjnskcjnkjnkjnkjnkjnkjn}{
    Sia $ F $ un campo di prevettori locale su $ \nonstandard{U} $ di classe $ D^{1} $. Se \begin{equation*}
        \norma{F(a)-F(b)-a+b}\le K \lambda\norma{a-b}
    \end{equation*}per ogni $ a,b \in \nonstandard{U}$, allora il flusso di $ F $ soddisfa: \begin{equation*}
        \norma{F^{n}(a)-F^{n}(b)-a+b}\le K \lambda ne^{K\lambda n}\norma{a-b}.
    \end{equation*}
}
\dimostrazionelem{djnckskdkjncjdkjnckjnskjndkjnckjnskcjnkjnkjnkjnkjnkjn}{
    Si ha che \begin{align*}
        \norma{F^{n}(a)-F^{n}(b)-a+b} &\le \sum_{i=1}^{n}\norma{F^{i}(a)-F^{i}(b)-F^{i-1}(a)+F^{i-1}(b)} \\ 
        &\le \sum_{i=1}^{n} K\lambda\norma{F^{i-1}(a)-F^{i-1}(b)}\\ 
        &\le \sum_{i=1}^{n} K \lambda(1+K\lambda)^{i-1}\norma{a-b} 
    \end{align*}dove l'ultima disuguaglianza è valida per \eqref{eq:cosa2}. Si ha inoltre che \begin{align*}
        \sum_{i=1}^{n}(1+K\lambda)^{i-1} &\le n(1+K\lambda)^{n-1}\le n e^{K \lambda n} 
    \end{align*}e questo dimostra la tesi.\qed
}
\proposizione{kdjnsckjndkjncskjndkjnckjnskjnckjnkjn}{
    Se $ F $ è un campo di prevettori di classe $ D^{1} $, allora $ \tilde{F} $ è di classe $ D^{1} $.}
\dimostrazioneprop{kdjnsckjndkjncskjndkjnckjnskjnckjnkjn}{
    Consideriamo una palla $ B $ racchiusa nel dominio di $ F $. Allora per ogni $ a,b \in \nonstandard{B}$ esiste $ K $ finito tale che \begin{equation*}
        \norma{F(a)-F(b)-a+b}\le K\lambda\norma{a-b}
    \end{equation*}per la Proposizione \framref{dlikjckjdlkmnckslkmdkcoslkdlkj}. Per $ n = \floor{t/\lambda} $ si ha, per il Lemma \framref{djnckskdkjncjdkjnckjnskjndkjnckjnskcjnkjnkjnkjnkjnkjn}, che \begin{equation*}
        \norma{F^{n}(a)-F^{n}(b)-a+b}\le K \lambda n e^{K\lambda n}\norma{a-b}\le K t e^{Kt}\norma{a-b}
    \end{equation*}Dunque si ha che $ \norma{h_{t}^{F} (a)-h_{t}^{F}(b)-a+b}\le K t e^{Kt}\norma{a-b} $ per ogni $ t $ reale e per ogni $ a,b $ standard, e, per il principio di transfer per $ t=\lambda $, \begin{equation*}
        \norma{\tilde{F}(a)-\tilde{F}(b)-a+b} \le  K \lambda e^{K\lambda}\norma{a-b}.\qedd
    \end{equation*}
}
\proposizione{kdjnsckjndkjncskjndkjnckjnskjnckjnkjnclassedidue}{
    Se $ F $ è un campo di prevettori di classe $ D^{2} $, allora $ \tilde{F} $ è di classe $ D^{2} $.}
\dimostrazioneprop{kdjnsckjndkjncskjndkjnckjnskjnckjnkjnclassedidue}{
    Per mezzo della Proposizione \framref{dkjncjdkjncjdkjncjdkjncjdjcjdkjncsdkjnkj} si assume che \begin{equation*}
        \norma{\Delta^{2}_{v,w}F(a)} \le K \lambda\norma{v}\norma{w}
    \end{equation*}per ogni $ a,v,w $ in qualche $ \nonstandard{B} $. Si dimostra ora per induzione interna che\footnote{Si ricorda che \begin{equation*}
        \norma{\Delta_{v,w}^{2}F^{n}(a) } = \norma{F^{n}(a)-F^{n}(a+v)-F^{n}(a+w)+F^{n}(a+v+w)}.
    \end{equation*}}\begin{equation*}
        \norma{\Delta_{v,w}^{2}F^{n}(a) } \le K \lambda\sum_{i=n-1}^{2n-2} (1+K\lambda)^{i}\norma{v}\norma{w}
    \end{equation*}per ogni $ n \in\nonstandard{\N} $.
    \begin{itemize}
        \item \emph{Passo base.} Il passo base per $ n=1 $ è l'ipotesi \begin{equation*}
            \norma{\Delta^{2}_{v,w}F(a)} \le K \lambda\norma{v}\norma{w}
        \end{equation*}
        \item \emph{Passo induttivo.} Si ipotizzi che l'affermazione valga per $ n $. Sia $ p=F^{n}(a) $, \begin{align*}
            x&=F^{n}(a+v)-F^{n}(a)\\
            y&=F^{n}(a+v)-F^{n}(a).
        \end{align*}Allora per la \eqref{eq:cosa2}:\begin{equation*}
            \norma{x}\le (1+K\lambda)^{n}\norma{v},\quad \norma{y}\le (1+K\lambda)^{n}\norma{w}
        \end{equation*}e, sfruttando l'ipotesi induttiva, \begin{align*}
            &\norma{F(p+x+y)- F^{n+1}(a+v+w)} \\
            &\hspace{4em} \le (1+K \lambda) \norma{p+x+y-F^n(a+v+w)}\\
            &\hspace{4em} =(1+K \lambda)\norma{-F^n(a)+F^n(a+v)+F^n(a+w)-F^n(a+v+w)}\\
            &\hspace{4em} \le (1+K \lambda) K \lambda \sum_{i=n-1}^{2 n-2}(1+K \lambda)^i\norma{v}\norma{w}\\
            &\hspace{4em}=K \lambda \sum_{i=n}^{2 n-1}(1+K \lambda)^i\norma{v}\norma{w},
            \end{align*}dove la prima disuguaglianza è per la Proposizione \framref{kjndjckjndjckjnskdkjncjdkskjncjdkjckjnjnjnjckjnkskjnc}, mentre la seconda è per ipotesi induttiva. 

            Ora si ha che \begin{align*}
                &\hspace{-3em}\norma{F^{n+1}(a)-F^{n+1}(a+v)-F^{n+1}(a+w)+F^{n+1}(a+v+w)}\\ 
                \le  &\norma{F(p)-F(p+x)-F(p+y)+F(p+x+y)}\\&+\norma{F(p+x+y)-F^{n+1}(a+v+w)}\\ 
                \le &K \lambda\norma{x}\norma{y} + K \lambda\sum_{i=n}^{2n-1}(1+K\lambda)^{i}\norma{v}\norma{w} \\ 
                \le &K \lambda(1+K\lambda)^{2n}\norma{v}\norma{w} + K \lambda\sum_{i=n}^{2n-1}(1+K\lambda)^{i}\norma{v}\norma{w} \\ 
                = & K \lambda\sum_{i=n}^{2n}(1+K\lambda)^{i}\norma{v}\norma{w}. 
            \end{align*}
    \end{itemize}

    Dunque per $ n=\floor{t/\lambda} $ si ha che \begin{multline*}
        \norma{F^{n}(a)-F^{n}(a+v)-F^{n}(a+w)+F^{n}(a+v+w)} \\
        \le K \lambda\sum_{i=n-1}^{2n-2}(1+K\lambda)^{i}\norma{v}\norma{w} \le K \lambda n e^{2Kt}\norma{v}\norma{w} \le K t e^{2Kt}\norma{v}\norma{w}.
    \end{multline*}Dunque, per standard $ a,v,w $ si ha che \begin{equation*}
        \norma{h_{t}^F(a)-h_{t}^F(a+v)-h_{t}^F(a+w)+h_{t}^F(a+v+w)}\le K t e^{2Kt}\norma{v}\norma{w},
    \end{equation*}e pertanto, estendendo al dominio non standard e valutando in $ t= \lambda $ si ottiene \begin{equation*}
        \norma{\tilde{F}(a)-\tilde{F}(a+v)-\tilde{F}(a+w)+\tilde{F}(a+v+w)}\le K t e^{2Kt}\norma{v}\norma{w}\qedd
    \end{equation*}
}
\lemma{kdjnckjnskjndjcjskjncjnsckjskjndjcjskjncjn}{
    Sia $ F $ un campo di prevettori locale su $ \nonstandard{U} $. Se $ \norma{F(a)-a} \le C\lambda$ e $ \norma{F(a)-a-F(b)+b}\le K \lambda\norma{a-b} $ per qualche $ C $ e $ K $ e per ogni $ a,b \in \nonstandard{U} $, allora il flusso di $ F $ soddisfa: \begin{equation*}
        \norma{F^{n}(a)-a-n\left(F(a)-a\right)} \le KCn^{2}\lambda^{2}.
    \end{equation*}
}
\dimostrazionelem{kdjnckjnskjndjcjskjncjnsckjskjndjcjskjncjn}{
    Si ha che \begin{align*}
        \norma{F^{n}(a)-a-n\left(F(a)-a\right)} &= \norma{\sum_{i=1}^{n} \left(F^{i}(a)-F^{i-1}(a)-\left(F(a)-a\right)\right) }\\ 
        &\le \sum_{i=1}^{n} \norma{F^{i}(a)-F^{i-1}(a)-\left(F(a)-a\right)} \\ 
        &= \sum_{i=1}^{n} \norma{F\left(F^{i-1}(a)\right)-F^{i-1}(a)-\left(F(a)-a\right)} \\&\le \sum_{i=1}^{n} K \lambda\norma{F^{i-1}(a)-a}\\ 
        &\le  \sum_{1\le j<i\le n}\norma{F^{j}(a)-F^{j-1}(a)} \le n^{2} K \lambda C\lambda\qedd 
    \end{align*}
}
\lemma{dkjnckjnsjdkjncjdkjnckjnkjnskjncjdjckjnsnncjdjdjjjjj}{
    Sia $ v \in \nonstandard{(\R^n)} $, con $ v\prec \lambda $, sia $ g:[0,+ \infty) \longrightarrow \R^{n} $ una funzione standard, \begin{equation*}
        g(t) = \st\left(\floor{t/\lambda}\,v\right)
    \end{equation*}e sia $ V= \st(v/\lambda). $

    Allora $ g(t)=tV $ e l'estensione di $ g $ al dominio non standard soddisfa $ g(\lambda)\equiv v $.
}
\dimostrazionelem{dkjnckjnsjdkjncjdkjnckjnkjnskjncjdjckjnsnncjdjdjjjjj}{
    Sia $ n=\floor{t/\lambda} $. Allora $ g(t) = \st(nv) $, e \begin{align*}
        \norma{tV-nv} &\le \norma{tV-t\left(\frac{v}{\lambda}\right)} + \norma{t\left(\frac{v}{\lambda}\right) - n\lambda\left(\frac{v}{\lambda}\right)}\\ 
        &= t\norma{V-v/\lambda} + \abs{t-n\lambda}\norma{v/\lambda} \pprec 1.
    \end{align*}Dunque $ g(t) = tV $, e quindi $ g(\lambda) = \lambda V $, e si ha che $ \norma{\lambda V-v} = \lambda\norma{V-v/\lambda}\pprec \lambda $.\qed
}
\teorema{dkjncjdjdkjnckjnskjndkjnckskjnckjndkjn}{
    Se $ F $ è un campo di prevettori locale di classe $ D^{1} $, allora $ \tilde{F}\equiv F $.
}
\dimostrazione{dkjncjdjdkjnckjnskjndkjnckskjnckjndkjn}{
    Se $ F $ è di classe $ D^1 $, allora $ \tilde{F} $ è di classe $ D^{1} $ per la Proposizione \framref{kdjnsckjndkjncskjndkjnckjnskjnckjnkjn}, e dunque è sufficiente mostrare che $ \tilde{F}(a)\equiv F(a) $ per ogni $ a $ standard, per il Lemma \framref{kjndkjncjdkjnckskdkjnjcsnkjn}.

    Dunque, per $ a $ standard e $ t $ standard, sia \begin{equation*}
        g_{t}(a) = \st\left(a+\floor{t/\lambda}\left(F(a)-a\right)\right) 
    \end{equation*}Denotando con $ n = \floor{t/\lambda} $ si ha \begin{align*}
        \norma{h_{t}^{F}-g_{t}(a)  } &= \norma{\st\left(F^n(a)\right)-\st\left(a+\floor{t/\lambda}\left(F(a)-a\right)\right)}\\ 
        &= \st\norma{F^{n}(a)-a-n\left(F(a)-a\right)}\le At^{2}
    \end{align*} per qualche $ A \in \R $, per il Lemma \framref{kdjnckjnskjndjcjskjncjnsckjskjndjcjskjncjn}. 
    
    Per il principio di transfer, valutando questo in $ t=\lambda $, si ottiene che \begin{equation*}
        \norma{\tilde{F}(a)-g_{\lambda}(a)} \le A\lambda^{2}\pprec \lambda
    \end{equation*}e dunque $ \tilde{F}(a)\equiv g_{\lambda}(a) $.

    Si ha che $ g_{t}(a)-a = \st\left(\floor{t/\lambda} \left(F(a)-a\right)\right) $ e dunque, per il Lemma \framref{dkjnckjnsjdkjncjdkjnckjnkjnskjncjdjckjnsnncjdjdjjjjj} $ g_{\lambda}(a)\equiv F(a)  $. Dal fatto che $ \equiv $ è una relazione di equivalenza segue la tesi.\qed
}
\section{Parentesi di Lie}
\definizione{
    Siano $ F,G $ campi di prevettori di classe $ D^{1} $. Si definiscono le \emph{parentesi di Lie} come segue: \begin{equation*}
        [F,G] = \left(G^{-1}\circ F^{-1}\circ G \circ F\right)^{\floor{1/\lambda}}.
    \end{equation*}
}

Si noti che $ F^{-1} $ e $ G^{-1} $ esistono per l'Osservazione \framref{remarkfourpointeight}.

\osservazione{
    Dal momento che fin dall'inizio $ \lambda $ è fissato e arbitrario, si può supporre che \begin{equation*}
        \lambda = \frac{1}{N},\quad N \in \nonstandard{\N}\setminus \N
    \end{equation*}e pertanto che $ 1/\lambda = N $ iperintero: in tal caso $ \floor{1/\lambda} = 1/\lambda $. D'ora in avanti, quindi, si scriverà $ 1/\lambda $ in luogo di $ \floor{1/\lambda} $, sulla base di questo assunto. 
}
\subsection{Proprietà fondamentali}
\teorema{dkjnckjnskjnkjnkjnkjnkknknkjnkkjnkkjnknkkjnkkjnkkjnk}{
    Se $ F,G $ sono campi di prevettori locali di classe $ D^{1} $, allora $ [F,G] $ è un campo di prevettori, i.e. $ [F,G](a)-a\prec \lambda $ per ogni $ a $.
}
\dimostrazione{dkjnckjnskjnkjnkjnkjnkknknkjnkkjnkkjnknkkjnkkjnkkjnk}{
    Poiché $ F $ è localmente di classe $ D^{1} $, vale che \begin{equation*}
        F(x)-x-F(y)+y \prec \lambda\norma{x-y}.
    \end{equation*}Sostituendo $ x=a $ e $ y=F^{-1}\circ G \circ F(a) $ si ottiene \begin{equation*}
        F(a)-a- G \circ F(a)+ F^{-1}\circ G\circ F(a) \prec \lambda\norma{a-F^{-1}\circ G\circ F(a)}\prec \lambda^{2}.
    \end{equation*}

    Poiché $ G $ è localmente di classe $ D^{1} $, vale che \begin{equation*}
        G(x)-x-G(y)+y \prec \lambda\norma{x-y}.
    \end{equation*}Sostituendo $ x=F(a) $ e $ y=G^{-1}\circ F^{-1}\circ G \circ F(a) $ si ottiene \begin{multline*}
        G\circ F(a)-F(a)- F^{-1}\circ G\circ F(a)+ G^{-1}\circ F^{-1}\circ G\circ F(a)\\ \prec \lambda\norma{F(a)-G^{-1}\circ F^{-1}\circ G\circ F(a)} \prec \lambda^{2}.
    \end{multline*}

    La somma delle due espressioni di cui sopra fornisce la seguente: \begin{equation*}
        G^{-1}\circ F^{-1}\circ G\circ F(a) -a \prec \lambda^{2}
    \end{equation*}Dunque, per \straniero{underspill}\footnote{Infatti, per $ \nonstandard{U} $ fissato, la relazione \begin{equation*}
        \norma{G^{-1}\circ F^{-1}\circ G\circ F(a)-a}\le C\lambda^{2}
    \end{equation*}vale per tutti i $ C\in \nonstandard{\N} $ infiniti.} in un appropriato $ \nonstandard{U}$, esiste $ C\prec 1 $ positivo tale che \[\norma{G^{-1}\circ F^{-1}\circ G\circ F(a)-a}\le C\lambda^{2}\] per ogni $ a \in \nonstandard{U} $. Pertanto: \begin{multline*}
        \norma{\left(G^{-1}\circ F^{-1}\circ G\circ F\right)^{1/\lambda}(a)-a}\\ 
        \le \sum_{k=1}^{1/\lambda} \norma{\left(G^{-1}\circ F^{-1}\circ G\circ F\right)^{k}(a)-\left(G^{-1}\circ F^{-1}\circ G\circ F\right)^{k-1}(a)}\le C\lambda.\qedd
    \end{multline*}
}
\lemma{kjndkckjnsdjckjndjdjcjdkjnccjdkjncjdjckjnsjdnj}{
    Siano $ F,G $ campi di prevettori locali, $ F $ di classe $ D^{2} $ e $ G $ di classe $ D^{1} $. Allora per ogni $ a,b $ con $ a-b\prec \lambda $ vale: \begin{equation*}
        F(a)-F(b)-F\circ G(a)+F\circ G(b) - a +b+G(a)-G(b)\prec \lambda^{2}\norma{a-b}.
    \end{equation*}
}
\dimostrazionelem{kjndkckjnsdjckjndjdjcjdkjnccjdkjncjdjckjnsjdnj}{
    Siano $ a,b $ fissati con $ a-b\prec \lambda $.

    Per la Proposizione \framref{dkjncjdkkjncsjdkjncsiujdoijfdosoidocodidislkdmlkmlkm} $ F $ è di classe $ D^{1} $, e pertanto per definizione \begin{equation*}
        F(\alpha)-\alpha-F(\beta) +\beta \prec \norma{\alpha-\beta}
    \end{equation*}con $ \alpha-\beta\prec \lambda $. Ponendo $ v=b-a $, $ w = G(a)-a $, $ \alpha = a+v+w $ e $ \beta = G(b) $, si ottiene \begin{equation*}
        F(a+v+w) - F\left(G(b)\right) - (a+v+w) + G(b)\prec \norma{a+v+w-G(b)}.
    \end{equation*}È banale verificare che $ \alpha,\beta $ così definiti soddisfino $ \alpha-\beta\prec\lambda $. 

    Inoltre, dal momento che $ a+v+w = b+G(a)-a $ si ha che \begin{equation*}
        F(a+v+w) - F\left(G(b)\right) - G(a)+a+G(b) \prec \norma{b+G(a)-a-G(b)}\prec \lambda^{2}\norma{a-b}
    \end{equation*}poiché $ G $ è di classe $ D^{1} $. Pertanto {\footnotesize\begin{align*}
        &\norma{F(a)-F(b)-F\circ G(a)+F\circ G(b) - a +b+G(a)-G(b)}\\ 
        &= \norma{F(a)-F(a+v)-F(a+w)+F(a+v+w)-F(a+v+w)+F\left(G(b)\right)-a+b+G(a)-G(b)}\\ 
        &\le  \norma{F(a)-F(a+v)-F(a+w)+F(a+v+w)}\\
        &\quad+\norma{-F(a+v+w)+F\left(G(b)\right)-a+b+G(a)-G(b)} \prec \lambda \norma{v}\norma{w} + \lambda^{2}\norma{a-b} \prec \lambda^{2}\norma{a-b}.\qedd
    \end{align*}}
}
\teorema{jkndkckdkjncjdjckjndjcjdkjncdkjncjdcj}{
    Se $ F,G $ sono campi di prevettori locali di classe $ D^{2} $, allora $ [F,G] $ è un campo di prevettori di classe $ D^{1} $.
}
\dimostrazione{jkndkckdkjncjdjckjndjcjdkjncdkjncjdcj}{
    Per la Proposizione \framref{dkjncjdkkjncsjdkjncsiujdoijfdosoidocodidislkdmlkmlkm}, la Proposizione \framref{kjndjckjnsdkjncjdjckjnsjdkcnskjncjdjdjcjskjnckskjnd} e per l'Osservazione \framref{dkcjnacjdjdkjcnkjnskjncjdjckjnsjcjdkjnc} si ha che entrambe \begin{equation*}
        F^{-1}\circ G\circ F,\qquad G^{-1}\circ F^{-1}\circ G
    \end{equation*}sono di classe $ D^{1} $.

    Si prenda ora in considerazione il Lemma \framref{kjndkckjnsdjckjndjdjcjdkjnccjdkjncjdjckjnsjdnj}: \begin{itemize}
        \item sostituendo nell'enunciato $ G $ con $ F^{-1}\circ G\circ F $ si ottiene, per $ a-b\prec \lambda $ \begin{equation*}
            F(a)-F(b)-G\circ F(a)+G\circ F(b)-a+b+F^{-1}\circ G\circ F(a)-F^{-1}\circ G\circ F(b)\prec \lambda^{2}\norma{a-b}
        \end{equation*}
        \item sostituendo nell'enunciato, rispettivamente, $ a,b,F,G $ con $ F(a),F(b), G, G^{-1}\circ F^{-1} $ si ottiene \begin{align*}
            &G\circ F(a)- G\circ F(b) - F^{-1}\circ G\circ F(a) + F^{-1}\circ G\circ F(b)\\ 
            &\quad - F(a)+F(b)+G^{-1}\circ F^{-1}\circ G\circ F(a) - G^{-1}\circ F^{-1}\circ G^{-1}\circ F(b)\\ 
            &\quad \prec \lambda^{2}\norma{F(a)-F(b)} \prec \lambda^{2}\norma{a-b}
        \end{align*}dove l'ultima relazione è dettata dalla Proposizione \framref{kjndjckjndjckjnskdkjncjdkskjncjdkjckjnjnjnjckjnkskjnc}. 
    \end{itemize}
    
    Sommando queste due disugualianze, si ottiene \begin{equation*}
        G^{-1}\circ F^{-1}\circ G\circ F(a)-G^{-1}\circ F^{-1}\circ G\circ F(b)-a+b \prec \lambda^2 \norma{a-b}.
    \end{equation*}Sia ora $ H\coloneqq G^{-1} \circ F^{-1}\circ G\circ F$: la relazione precedente si riscrive come \[H(a) - H(b) -a+b \prec \lambda^{2}\norma{a-b};\]inoltre si ha che $ [F, G] = H^{1/\lambda} $. Per \straniero{underspill} su un appropriato $ \nonstandard{U} $, esiste $ C\prec 1 $ tale per cui \begin{equation*}
        \norma{H(a)-H(b)-a+b}\le C\lambda^{2}\norma{a-b}
    \end{equation*}per ogni $ a,b \in \nonstandard{U} $. Allora, per il Lemma \framref{djnckskdkjncjdkjnckjnskjndkjnckjnskcjnkjnkjnkjnkjnkjn}, posto $ K=C\lambda $ e $ n=1/\lambda $, si ottiene \begin{equation*}
        \norma{H^{1/\lambda}(a)-H^{1/\lambda}(b)-a+b}\le C\lambda e^{C\lambda}\norma{a-b}.\qedd
    \end{equation*}
}
\proposizione{kjndkjnckjnskdjnckjnskjdjcjdkjncjdjckjncdkjnckjnckjns}{
    Siano $ F,G $ campi di prevettori di classe $ D^{2}$. Allora, per ogni $ a $, \begin{equation*}
        [F,G](a)\equiv a+ \frac{1}{\lambda}\left(G^{-1}\circ F^{-1}\circ G\circ F(a)-a\right)
    \end{equation*}
}
\dimostrazioneprop{kjndkjnckjnskdjnckjnskjdjcjdkjncjdjckjncdkjnckjnckjns}{
    Sia $ H\coloneqq G^{-1} \circ F^{-1}\circ G\circ F$. La dimostrazione del Teorema \teoref{dkjnckjnskjnkjnkjnkjnkknknkjnkkjnkkjnknkkjnkkjnkkjnk} garantisce che esista $ C'\prec 1 $ tale per cui $ \norma{H(a)-a}\le C' \lambda^{2} $ per ogni $ a $.

    La dimostrazione del Teorema \teoref{jkndkckdkjncjdjckjndjcjdkjncdkjncjdcj} garantisce che esista $ C''\prec 1 $ tale per cui \[\norma{H(a)-H(b)-a+b}\le C''\lambda^{2}\norma{a-b}\]per ogni $ a,b $.
    
    Dunque, considerando il Lemma \framref{djnckskdkjncjdkjnckjnskjndkjnckjnskcjnkjnkjnkjnkjnkjn}, e ponendo $ C = C'\lambda $ e $ K = C'' \lambda $ e $ n=1/\lambda $, si ottiene che \begin{equation*}
        \norma{H^{1/\lambda}(a)-a - \frac{1}{\lambda} \left(H(a)-a\right)} \le C'' \lambda \cdot C' \lambda \cdot \frac{1}{\lambda^{2}} \lambda^{2} = C'' C'' \lambda^{2}\pprec \lambda\qedd 
    \end{equation*}
}
\lemma{dkjnckjnskjdjckjncjdjdjdjdjlemminoinoinoinoinoinoinoinoi}{
    Se $ G,H $ sono campi di prevettori di classe $ D^{2} $ e $ G\equiv H $, allora \begin{equation*}
        \left(G(a)-H(a)\right)-\left(G(b)-H(b)\right)\pprec \lambda\norma{a-b}
    \end{equation*}per ogni $ a,b $ con $ a-b\prec \lambda $.
}
\dimostrazionelem{dkjnckjnskjdjckjncjdjdjdjdjlemminoinoinoinoinoinoinoinoi}{
    Sia $ F(x) = G(x)-H(x) $: si ha che $ F(x) \pprec \lambda $ per ogni $ x $ (questo è garantito dall'ipotesi $ G\equiv H $). 

    Se per assurdo $ F(a)-F(b)$ non è $ \pprec \lambda\norma{a-b} $ per qualche $ a,b $, con $ a-b\prec \lambda $, allora \begin{equation*}
        \lambda\norma{a-b}\prec \norma{F(a)-F(b)}.
    \end{equation*}Sia $ v = b-a $, allora \begin{equation*}
        \lambda\norma{v}^{2}\prec \norma{v}\norma{F(a)-F(a+v)},
    \end{equation*}e, poiché $ G $ e $ H $ sono $ D^{2} $, $ \Delta^{2}_{v,v} F(x) \prec \lambda\norma{v}^{2} $ per ogni $ x $. Dunque per ogni $ x $\begin{equation*}
        \Delta_{v,v}^{2}F(x) \prec \norma{v}\,\norma{F(a)-F(a+v)} .
    \end{equation*}Per il Lemma \framref{djjjjjjjjjframmentinoframmentucciovieniqui}, prendendo una palla aperta attorno ad $ \st(a) $, esiste $ m \in \nonstandard{N} $ tale per cui \begin{align*}
        F(a)-F(a+v)&\prec \norma{v}\norma{F(a) - F(a+mv)}\\ 
        & \le \norma{v}\left(\norma{F(a)}+\norma{F(a+mv)}\right)\pprec \norma{v}\lambda.\qedd
    \end{align*}
}
\lemma{dlknckjnskjnckjdjjjjjdjjjdlemininiosjndksjnckjn}{
    Se $ G,H $ sono campi di prevettori con $ G\equiv H $, e $ G $ è di classe $ D^{1} $, e $ H^{-1} $ esiste, allora $ G^{-1}\equiv H^{-1} $.
}
\dimostrazionelem{dlknckjnskjnckjdjjjjjdjjjdlemininiosjndksjnckjn}{
    Fissato $ a $, siano \begin{equation*}
        x= G^{-1}(a),\quad y = H^{-1}(a);
    \end{equation*}si deve mostrare che $ x-y\pprec \lambda $. 

    Poichè $ G(x) = a = H(y) $, vale che \begin{equation*}
        \norma{G(x) - G(y)} = \norma{H(y)-G(y)} = \beta\lambda
    \end{equation*}per qualche $ \beta \pprec 1 $ (poiché $ G\equiv H $). Poiché $ G $ è di classe $ D^{1} $:\[\norma{G(x)-G(y)-x+y} = K \lambda\norma{x-y}\]per qualche $ K \prec 1 $, e pertanto \begin{equation*}
        \norma{x-y} \le \norma{G(x)-G(y)-x+y} + \norma{G(x)-G(y)} = K\lambda\norma{x-y}+\beta\lambda
    \end{equation*}e, nello specifico $ (1-K\lambda)\norma{x-y} \le \beta\lambda $. Questo dimostra la tesi, in quanto \begin{equation*}
        \norma{x-y}\le \frac{\beta}{1-K\lambda}\lambda\pprec \lambda.\qedd
    \end{equation*}
}
\teorema{dkjncjdkjnckjnskdjjjjjdjjdkjncjdkkjncjskdlkcsk}{
    Se $ F,E,G,H $ sono campi di prevettori di classe $ D^{2} $ e $ F\equiv E $, $ G\equiv H $, allora $ [F,G]\equiv [E,H] $.
}
\dimostrazione{dkjncjdkjnckjnskdjjjjjdjjdkjncjdkkjncjskdlkcsk}{
    È sufficiente mostrare il caso per $ F=E $. Infatti, dalla definizione è chiaro che $ [G,F]=[F,G]^{-1} $, e dunque, se $ [F,G]\equiv [F,H] $ e $ [H,F]\equiv [H,E] $, allora \begin{equation*}
        [F,G]\equiv [F,H] = [H,F]^{-1} \equiv [H, E]^{-1} = [E,H]
    \end{equation*}per il Teorema \teoref{jkndkckdkjncjdjckjndjcjdkjncdkjncjdcj} e il Lemma \framref{dlknckjnskjnckjdjjjjjdjjjdlemininiosjndksjnckjn}. 

    Si procede dunque a dimostrare che $ [F,G]\equiv [F,H] $. 

    Fissato $ x $, siano \begin{equation*}
        a= G\left(F(x)\right), \quad b = H\left(F(x)\right);
    \end{equation*}per ipotesi, $ a-b\pprec \lambda $. Per la Proposizione \framref{dkjncjdkkjncsjdkjncsiujdoijfdosoidocodidislkdmlkmlkm} e la Proposizione \framref{kjndjckjnsdkjncjdjckjnsjdkcnskjncjdjdjcjskjnckskjnd} \begin{equation*}
        F^{-1}(a)-F^{-1}(b)-a+b\prec \lambda\norma{a-b}\pprec \lambda^{2}.
    \end{equation*}Sia ora $ c \coloneqq G^{-1}\circ F^{-1}\circ G\circ F(x) $. Per il Lemma \framref{dkjnckjnskjdjckjncjdjdjdjdjlemminoinoinoinoinoinoinoinoi} \begin{equation*}
        F^{-1}(a)-H(c) - a +b = \left(G(c)-H(c)\right) - \left(G\left(F(x)\right)-H\left(F(x)\right)\right) \pprec \lambda \norma{c-F(x)}\pprec \lambda^{2}
    \end{equation*}Sommando le ultime due espressioni si ottiene che $ H(c) - F^{-1}(b)\pprec \lambda^{2} $ e pertanto \begin{multline*}
        G^{-1}\circ F^{-1}\circ G\circ F(x) - H^{-1}\circ F^{-1}\circ H\circ F(x) =\\ 
        = H^{-1}\left(H(c)\right) - H^{1}\left(F^{-1}(b)\right)\prec \norma{H(c)-F^{-1}(b)}\pprec \lambda^{3}
    \end{multline*}per la Proposizione \framref{dkjncjdkkjncsjdkjncsiujdoijfdosoidocodidislkdmlkmlkm}, la Proposizione \framref{kjndjckjnsdkjncjdjckjnsjdkcnskjncjdjdjcjskjnckskjnd} e la Proposizione \framref{kjndjckjndjckjnskdkjncjdkskjncjdkjckjnjnjnjckjnkskjnc}. Dunque vale \begin{equation*}
        \frac{1}{\lambda}\left(G^{-1}\circ F^{-1}\circ G\circ F(x)\right) - \frac{1}{\lambda}\left(H^{-1}\circ F^{-1}\circ H\circ F(x)\right) \pprec \lambda
    \end{equation*}e dunque, per la Proposizione \framref{kjndkjnckjnskdjnckjnskjdjcjdkjncjdjckjncdkjnckjnckjns} $ [F,G](x)\equiv [F,H](x) $.\qed
}
\subsection{Relazione con la parentesi di Lie classica}

\teorema{dkjnckskjndjjjdjjdjjdkjnckjnskjnkjnkjnkjnkjnkjnkjnkkjnkjnkjnkjkjnkjnjnjjnjnj}{
    Siano $ X,Y $ due campi vettoriali di classe $ C^{2} $, e sia $ [X,Y]_{\text{cl}} $ le loro parentesi di Lie classiche. 

    Siano $ F,G $ due campi di prevettori che realizzino $ X,Y $ rispettivamente. Allora $ [F,G] $ realizza $ [X,Y]_{\text{cl}} $.
}
\dimostrazione{dkjnckskjndjjjdjjdjjdkjnckjnskjnkjnkjnkjnkjnkjnkjnkkjnkjnkjnkjkjnkjnjnjjnjnj}{
    Per l'Osservazione \framref{osservsADJKNCkjnadkjhdbcakjfhakjdhfgaskjdfghadsfjgh}, i flussi $ h^{F}_t, h^{G}_t $ coincidono con i flussi di $ X,Y $ nel senso classico\footnote{È possibile sfruttare l'osservazione citata grazie al Teorema \teoref{dkjncjdkjnckjnskdjjjjjdjjdkjncjdkkjncjskdlkcsk}.}. È ben noto che, in relazione al flusso classico, \begin{equation*}
        [X,Y]_{\text{cl}} (p) = \lim_{t\to 0} \frac{1}{t^{2}}\left(\left(h_{t}^{G} \right)^{-1} \circ \left(h_{t}^{F} \right)^{-1}\circ h_{t}^{G}\circ h_{t}^{F} - p \right).
    \end{equation*}Per la caratterizzazione dei limiti attraverso gli infinitesimi, si ha \begin{equation*}
        [X,Y]_{\text{cl}} (p) \approx \frac{1}{\lambda^{2}}\left(\tilde{G}^{-1}\circ\tilde{F}^{-1}\circ \tilde{G}\circ \tilde{F}(p)-p\right).
    \end{equation*}Se $ v\approx w $, allora $ \lambda v\equiv \lambda w $ e pertanto, per la procedura introdotta nella dimostrazione della Proposizione \framref{dkjnckjdjdjdjjjjjdjjjdicoij}, $ [X,Y]_{\text{cl}} $ è realizzato dal campo di prevettori \begin{equation*}
        A(a) = a + \frac{1}{\lambda}\left(\tilde{G}^{-1}\circ\tilde{F}^{-1}\circ \tilde{G}\circ \tilde{F}(p)-p\right).
    \end{equation*}

    Resta da dimostrare che $ [F,G]\equiv A $. Per la Proposizione \framref{kdjnsckjndkjncskjndkjnckjnskjnckjnkjnclassedidue}, $ \tilde{F},\tilde{G} $ sono di classe $ D^{2} $, e pertanto per la Proposizione \framref{kjndkjnckjnskdjnckjnskjdjcjdkjncjdjckjncdkjnckjnckjns}, $ [\tilde{F},\tilde{G}]\equiv A $. Per il Teorema \teoref{dkjncjdjdkjnckjnskjndkjnckskjnckjndkjn} (tramite la Proposizione \framref{dkjncjdkkjncsjdkjncsiujdoijfdosoidocodidislkdmlkmlkm}), $ F\equiv \tilde{F} $, $ G\equiv \tilde{G} $, e pertanto, per il Teorema \teoref{dkjncjdkjnckjnskdjjjjjdjjdkjncjdkkjncjskdlkcsk}, $ [F,G]\equiv A $.\qed
}

\section{Conclusione}
Si conclude quindi con quello che è l'obiettivo di questa tesi: un corollario classico, dimostrato con gli strumenti dell'analisi non standard.

Si dimostra prima la versione per la parentesi di Lie non standard, e poi, come corollario, tramite il Teorema \teoref{dkjnckskjndjjjdjjdjjdkjnckjnskjnkjnkjnkjnkjnkjnkjnkkjnkjnkjnkjkjnkjnjnjjnjnj}, si dimostra per la parentesi di Lie standard. 

\teorema{Skjndkjnckjnsjdkjncjndkjncjdkjnsdjckjn}{
    Siano $ F,G $ due campi di prevettori di classe $ D^{2} $. 

    Allora $ [F,G]\equiv I $ se e solo se esiste $ 0<T \in \R $ tale che per ogni $ 0\le t,s\le T $:\[h_{t}^{F}\circ h_{s}^{G} = h_{s}^{G}\circ h_{t}^{F}.\]
}
\dimostrazione{Skjndkjnckjnsjdkjncjndkjncjdkjnsdjckjn}{
    Si assuma che $ [F,G]\equiv I $, i.e. \begin{equation*}
        [F,G](a)-a\pprec \lambda
    \end{equation*}per ogni $ a $. Per la Proposizione \framref{kjndkjnckjnskdjnckjnskjdjcjdkjncjdjckjncdkjnckjnckjns} 
    \begin{equation*}
        \frac{1}{\lambda}\left(G^{-1}\circ F^{-1}\circ G\circ F ( a) - a\right)\pprec \lambda
    \end{equation*}e dunque $ \left(G^{-1}\circ F^{-1}\circ G\circ F ( a) - a\right)\pprec \lambda^{2} $ per ogni $ a $. Questo implica, per la Proposizione \framref{kjndjckjndjckjnskdkjncjdkskjncjdkjckjnjnjnjckjnkskjnc} che \begin{equation*}
        G\circ F(a)-F\circ G(a)\pprec \lambda^{2}\qquad \forall\, a.
    \end{equation*}

    Siano quindi $ n = \floor{t/\lambda}, m =\floor{s/\lambda} $: è necessario mostrare che $ F^{n}\circ G^{m}(a)\approx G^{m}\circ F^{n}(a) $ per ogni $ a $. Questo processo comporta $ nm $ scambi di $ F $ e $ G $; ciascuno scambio è nella forma \begin{equation*}
        F^{k}\circ G^{r}\circ F \circ G^{m-r}\circ F^{n-k-1} \quad \longrightarrow \quad F^{k}\circ G^{r+1}\circ F\circ G^{m-r-1} \circ F^{n-k-1}
    \end{equation*}Applicando $ G\circ F(p)-F\circ G(p)\pprec \lambda^{2} $ a $ p = G^{m-r-1}\circ F^{n-k-1}(a) $ si ottiene \begin{equation*}
        F\circ G^{m-r} \circ F^{n-k-1}(a)-G\circ F\circ G^{m-r-1}\circ F^{n-k-1}(a)\pprec \lambda^{2}
    \end{equation*}

    Per la Proposizione \framref{dkjncjdkkjncsjdkjncsiujdoijfdosoidocodidislkdmlkmlkm} e per la Proposizione \framref{dlikjckjdlkmnckslkmdkcoslkdlkj} esiste $ K \in \R $ tale per cui \begin{align*}
        \norma{F(a)-a-F(b)+b} &\le K \lambda\norma{a-b}\\ 
        \norma{G(a)-a-G(b)+b} &\le K \lambda\norma{a-b}
    \end{align*}per ogni $ a,b $ in un dominio appropriato. 
    Allora, per il Teorema \teoref{dlkmncldlslkmckdlclkmsdkckdkdkkkkkdkk} applicato a $ G^{r} $ e poi a $ F^{k} $ si ottiene che \begin{multline*}
        \norma{F^{k}\circ G^{r} \circ F\circ G^{m-r} \circ F^{n-k-1}(a)- F^{k}\circ G^{r+1} \circ F\circ G^{m-r-1}\circ F^{n-k-1}(a)}\\ 
        \le e^{K(t+s)}\norma{F\circ G^{m-r} \circ F^{n-k-1}(a)-G\circ F\circ G^{m-r-1}\circ F^{n-k-1}(a)}\pprec \lambda^{2}
    \end{multline*}Dunque, sommando gli $ nm$ contributi passando da $ F^{n}\circ G^{m}(a) $ a $ G^{m}\circ F^{n}(a)$ si ottiene che \begin{equation*}
        F^{n}\circ G^{m}(a)-G^{m}\circ F^{n}(a)\pprec 1.
    \end{equation*}Questo perché tra le $ nm $ differenze che si sommano, ce ne sarà una massimale (sia ad esempio $ \beta \lambda^{2} $), e pertanto la somma di tutti gli $ nm $ contributi sarà $ \le nm\beta \lambda^{2}\le ts\beta \pprec 1 $. Si è dimostrata la prima implicazione.

    Viceversa, se $ h_{t}^{F}\circ h_{t}^{G} = h_{t}^{G}\circ h_{t}^{F}  $. Allora per il principio di transfer, $ \tilde{F}\circ \tilde{G} = \tilde{G} \circ \tilde{F} $. e pertanto \begin{equation*}
        \tilde{G}^{-1}\circ \tilde{F}^{-1} \circ \tilde{G}\circ\tilde{F} = I
    \end{equation*}e dunque $ [\tilde{F},\tilde{G}]=I $. Per la Proposizione \framref{kdjnsckjndkjncskjndkjnckjnskjnckjnkjnclassedidue} e per il Teorema \teoref{dkjncjdjdkjnckjnskjndkjnckskjnckjndkjn}, e il Teorema \teoref{dkjncjdkjnckjnskdjjjjjdjjdkjncjdkkjncjskdlkcsk}  si ottiene \[[F,G]\equiv I .\qedd\]
}

\teorema[Corollario]{djnckjnsjdjckjnskjndjcjdkjncjdkjncjdjdkjnc}{
    Siano $ X,Y $ dei campi vettoriali classici di classe $ C^{2} $. 
    
    Allora i flussi di $ X $ e di $ Y $ commutano se e solo se $ [X,Y]_{\text{cl}}=0 $.
}
\dimostrazione{djnckjnsjdjckjnskjndjcjdkjncjdkjncjdjdkjnc}{
    Siano $ F,G $ campi di prevettori definiti come nella dimostrazione della Proposizione \framref{dkjnckjdjdjdjjjjjdjjjdicoij}, ovvero \begin{align*}
        F(a) &= a + \lambda X(a)\\ 
        G(a) &= a+ \lambda Y(a).
    \end{align*}Applicando la Proposizione \framref{dcoickjdkjnckjnkskjnksckjndkjnckjnckjdkjnckkjnksjncj} si ottiene che $ F,G $ sono di classe $ D^{2} $, e per l'Osservazione \framref{osservsADJKNCkjnadkjhdbcakjfhakjdhfgaskjdfghadsfjgh} $ h_{t}^{F}  $ e $ h_{t}^{G}  $ rappresentano i flussi di, rispettivamente, $ X $ e $ Y $ nel senso classico. 

    La tesi è ora banale applicazione del Teorema \teoref{dkjnckskjndjjjdjjdjjdkjnckjnskjnkjnkjnkjnkjnkjnkjnkkjnkjnkjnkjkjnkjnjnjjnjnj} e del Teorema \teoref{Skjndkjnckjnsjdkjncjndkjncjdkjnsdjckjn}, tenendo a mente l'Osservazione \framref{kjndkjncskjndjjjdjjdjdjdjjdsjhbcsjhdbcs}.\qed
}
